%% 8-7 道にそってたす：線積分（第14回：組 V8-4。ページ制限なし。偶数ページ）
%% 「名前と操作」の型（tools/kata8-int.md）で書く。場合分けの表8-7（A〜K）と流れ図。
%% ほかの巻・物理の本との約束：物理 p1-04（力-4）の「道を小さな移動に分け…区切りを細かくした先が線積分」、
%%   F1=(y,x)（U=−xy）・F2=(−y,x)、道A・道B（F2 は 1 と −1、F1 はどちらも1）、円を100本の線分で 6.2791、
%%   「すきまのない広がりの中なら」（8-10）、「rot の値2×囲んだ面積1」（8-11）。4-12（内積・仕事）、5-12、
%%   主役の山を登る仕事 8330 J。p-V 図 400・200・300 J（熱-2 の予告、S12）。変分法の入口（S1）はもっと知りたい。
%%   数はすべて Python で検算した（scratchpad/v8-4/check.py）。
\makeatletter
%% ===== 「名前と操作」の型（tools/kata8-int.md、組 V8-4）。頭文字 kviiih は章ごとに置きかえる =====
\colorlet{kviiihNc}{bkOrange!85!black}% 名前の枠の線（点線）
\colorlet{kviiihSc}{bkTheme}% ふつうの定積分の枠の線（実線）
\colorlet{kviiihSb}{bkTheme!8}% ふつうの定積分の地
\tikzset{
  kviiih dan/.style={font=\sffamily\bfseries\footnotesize, rounded corners=1.2mm, inner xsep=1.6mm, inner ysep=0.5mm, line width=0.9pt},
  kviiih dan1/.style={kviiih dan, draw=kviiihNc, dash pattern=on 2.2pt off 1.4pt, fill=white, text=bkInk},
  kviiih dan2/.style={kviiih dan, draw=kviiihSc, fill=kviiihSb, text=bkInk},
  kviiih dan3/.style={kviiih dan, draw=bkTheme, fill=bkTheme, text=bkOn},
  kviiih name/.style={draw=kviiihNc, line width=0.9pt, dash pattern=on 2.2pt off 1.4pt, rounded corners=1.2mm, inner xsep=1.2mm, inner ysep=1.2mm, fill=white},
  kviiih seki/.style={draw=kviiihSc, line width=0.9pt, rounded corners=1.2mm, inner xsep=1.2mm, inner ysep=1.2mm, fill=kviiihSb},
  kviiih tag/.style={font=\sffamily\bfseries\footnotesize, text=bkInk, inner sep=0.4mm},
  kviiih q/.style={diamond, aspect=2.4, draw=bkYellow!75!black, line width=0.9pt, fill=bkYellowBg, font=\sffamily\footnotesize, align=center, inner xsep=1.2mm, inner ysep=0.8mm, text=bkInk},
  kviiih do/.style={rectangle, draw=\bkTT, line width=0.9pt, fill=white, rounded corners=0.6mm, font=\sffamily\footnotesize, align=center, inner xsep=1.5mm, inner ysep=1.2mm, text=bkInk},
  % 流れ図（Ver2：矢じり 2.6 mm、先は枠から 0.5 mm。「はい」は矢印の上、「いいえ」は矢印の右。白地でかくさない）
  kviiih go/.style={-{Stealth[length=2.6mm, width=2.2mm]}, line width=0.9pt, bkInk!75, shorten <=0.5mm, shorten >=0.5mm},
  kviiih yn/.style={font=\sffamily\bfseries\footnotesize, text=bkInk, inner sep=0.5mm},
  kviiih yes/.style={kviiih yn, midway, above=0.3mm},
  kviiih no/.style={kviiih yn, midway, right=0.6mm},
  kviiih note/.style={font=\sffamily\footnotesize, text=bkInk, align=left, text width=132mm, inner sep=0pt},
}
% 段の札（例題の中で）  \kviiihdan{1|2|3}
\newcommand\kviiihdan[1]{\tikz[baseline=(d.base)]\node[kviiih dan#1] (d) {\vphantom{あ}\ifcase#1\or 名前\or 直す\or 楽をする\fi};}
% 段の見出し（本文で）  \kviiihdanhead{1|2|3}{見出し}
\newcommand\kviiihdanhead[2]{\par\addvspace{8pt}\Needspace{7\baselineskip}\noindent
  \tikz[baseline=(d.base)]\node[kviiih dan#1, font=\sffamily\bfseries\small, inner ysep=0.8mm] (d) {\vphantom{あ}\ifcase#1\or ①\,名前の段\or ②\,直す段\or ③\,楽をする段\fi};%
  \hspace{0.7\zw}{\sffamily\bfseries\color{\bkTT}#2}\par\nobreak\vspace{3pt}\noindent\ignorespaces}
% 場合の札  \kviiihcase{C}
\newcommand\kviiihcase[1]{\tikz[baseline=(c.base)]\node[fill=bkTheme, text=bkOn, font=\sffamily\bfseries\small, rounded corners=0.8mm, inner xsep=1mm, inner ysep=0.6mm, minimum width=5.4mm] (c) {\vphantom{Ag}#1};}
% 例題の1行め  \kviiihbaai{ことば}{8-6}{C}（場合が2つなら {C,G}）
\newcommand\kviiihbaai[3]{{\sffamily\bfseries 場合：}#1（表#2の\kviiih@cases#3,\@nil）}
\def\kviiih@cases#1,#2\@nil{\kviiihcase{#1}\ifx\relax#2\relax\else と\kviiih@cases#2\@nil\fi}
% 式の中の「名前」と「ふつうの定積分」
\newcommand\kviiihN[1]{\tikz[baseline=(n.base)]\node[kviiih name] (n) {$\displaystyle #1$};}
\newcommand\kviiihS[1]{\tikz[baseline=(n.base)]\node[kviiih seki] (n) {$\displaystyle #1$};}
\newcommand\kviiihNt[1]{\tikz[baseline=(n.base)]{\node[kviiih name] (n) {$\displaystyle #1$};\node[kviiih tag, anchor=north] at ([yshift=-0.8mm]n.south) {名前};}}
\newcommand\kviiihSt[1]{\tikz[baseline=(n.base)]{\node[kviiih seki] (n) {$\displaystyle #1$};\node[kviiih tag, anchor=north] at ([yshift=-0.8mm]n.south) {ふつうの定積分};}}
% 場合分けの表の題の帯  \kviiihcasetitle{8-6}{題}
\newcommand\kviiihcasetitle[2]{\par\addvspace{6pt}\Needspace{8\baselineskip}\noindent\begin{tikzpicture}[baseline=(t.base)]
  \node[fill=bkTheme, text=bkOn, font=\sffamily\bfseries\small, rounded corners=1mm, inner xsep=2mm, inner ysep=1mm] (t) {場合分けの表\,#1};
  \node[anchor=west, font=\sffamily\bfseries, text=\bkTT, inner sep=0pt] at ([xshift=2mm]t.east) {#2};
  \end{tikzpicture}\par\nobreak\vspace{3pt}}
% タイルの小さな絵（対応表の左の列）  \kviiihpic{1〜7}
\newcommand\kviiihpic[1]{\begin{tikzpicture}[x=1mm, y=1mm, baseline=(current bounding box.center)]
  \ifcase#1\or
  \draw[bkInk, line width=0.6pt] (0,0) -- (14,0);
  \foreach \kviiihx/\kviiihh in {3/3.9,5/4.6,7/4.6,9/4.5,11/5.1} {\filldraw[fill=bkTheme!25, draw=bkTheme, line width=0.4pt] ({\kviiihx-1},0) rectangle ({\kviiihx+1},\kviiihh);}
  \draw[bkInk!80, line width=0.7pt] (1,3.2) .. controls (5,6) and (9,3) .. (13,5.6);
  \or
  \draw[bkGray, line width=0.4pt, step=2.5] (0,0) grid (12.5,7.5);
  \fill[bkTheme!45] (5,2.5) rectangle (7.5,5);
  \draw[bkInk, line width=0.7pt] (0,0) rectangle (12.5,7.5);
  \or
  \draw[bkInk, line width=0.6pt] (0,0) rectangle (9,6) (0,6) -- (3,8.5) -- (12,8.5) -- (9,6) (12,8.5) -- (12,2.5) -- (9,0);
  \filldraw[fill=bkTheme!45, draw=bkTheme, line width=0.4pt] (3.5,1.5) rectangle (5.5,3.5) (3.5,3.5) -- (4.3,4.2) -- (6.3,4.2) -- (5.5,3.5) (6.3,4.2) -- (6.3,2.2) -- (5.5,1.5);
  \or
  \draw[bkInk, line width=0.7pt] (0,1) .. controls (4,8) and (8,-1) .. (13,6);
  \draw[bkTheme, line width=1.8pt] (5.2,3.75) -- (7.4,2.9);
  \or
  \draw[bkInk, line width=0.7pt] (0,1) .. controls (4,8) and (8,-1) .. (13,6);
  \draw[bkTheme, line width=1.3pt, -{Stealth[length=1.8mm]}] (4.6,3.95) -- (7.8,2.85);
  \or
  \draw[bkInk, line width=0.6pt] (0,1) .. controls (4,4) and (8,4) .. (12,1) -- (13,6) .. controls (9,9) and (5,9) .. (1,6) -- cycle;
  \filldraw[fill=bkTheme!45, draw=bkTheme, line width=0.4pt] (5,3.6) -- (7.4,3.6) -- (7.8,5.8) -- (5.4,5.8) -- cycle;
  \or
  \draw[bkInk, line width=0.6pt] (0,1) .. controls (4,4) and (8,4) .. (12,1) -- (13,6) .. controls (9,9) and (5,9) .. (1,6) -- cycle;
  \filldraw[fill=bkTheme!45, draw=bkTheme, line width=0.4pt] (5,3.6) -- (7.4,3.6) -- (7.8,5.8) -- (5.4,5.8) -- cycle;
  \draw[bkRed, line width=1.1pt, -{Stealth[length=1.8mm]}] (6.4,4.7) -- (6.4,10.5);
  \fi
\end{tikzpicture}}
%% ===== この章の部品 =====
% 小見出し（太字の1行）
\newcommand\kviiihhead[1]{\par\addvspace{6pt}\Needspace{6\baselineskip}\noindent{\sffamily\bfseries\color{\bkTT}#1}\par\nobreak\vspace{2pt}\noindent\ignorespaces}
% 囲みの中の小見出し（それだけで1行）
\newcommand\kviiihbox[1]{\par\addvspace{3pt}\noindent\textbf{#1}\par\vspace{3pt}\noindent\ignorespaces}
% 式を、まん中に1行で置く
\newcommand\kviiihline[1]{\par\addvspace{3pt}{\centering #1\par}\addvspace{3pt}\noindent\ignorespaces}
% 大事な式のわく  \kviiihframe[はば]{中身}
\newcommand\kviiihframe[2][100mm]{\par\addvspace{4pt}{\centering\tikz\node[draw=\bkTT, line width=0.9pt, fill=bkTheme!6, rounded corners=1.5mm, inner xsep=6mm, inner ysep=2.5mm, align=center]{#2};\par}\addvspace{4pt}\noindent\ignorespaces}
% 例題の途中式の2行め以降
\newcommand\kviiiha{\par\vspace{2pt}\noindent\rightskip=0pt plus 8em\relax\hangindent=2.6\zw\hangafter=0\relax}
\newcommand\kviiihb{\par\vspace{2pt}\noindent\rightskip=0pt plus 8em\relax\hangindent=4.3\zw\hangafter=0\relax}
\newcommand\kviiihs{\rightskip=0pt plus 8em\relax}
\makeatother
%% ===== この章の部品（8-7）=====
\makeatletter
% 平面の矢印の場（F2=(−y, x) など）を格子点にかく  \kviiihfield{Fx の式}{Fy の式}{1目もり(mm)}{矢印の倍率(mm)}{最大}
\newcommand\kviiihfield[5]{\foreach \kviiihi in {0,0.5,...,#5} {\foreach \kviiihj in {0,0.5,...,#5} {%
  \pgfmathsetmacro\kviiihu{#1}\pgfmathsetmacro\kviiihv{#2}%
  \draw[bkGray!80!black, line width=0.7pt, -{Stealth[length=1.6mm]}] ({\kviiihi*#3},{\kviiihj*#3}) -- ++({\kviiihu*#4},{\kviiihv*#4});}}}
\makeatother

\bkchapter{8}{7}{道にそってたす：\\線積分}{\bklevel{5}}{道にそってたす：線積分}
\bkmichishirube{4-11,4-12,5-8,5-12,8-3,8-4,8-5,8-6}{8-8,8-10,8-11,8-14}

\begin{goal}
\begin{itemize}[label={\color{\bkTT}●}, leftmargin=1.6\zw, itemsep=0pt, topsep=1pt]
\item \bklevel{3}\enspace 道を小さな1歩に分け、「力と1歩の内積」をたして仕事を表計算で求め、線積分の記号を読める。
\item \bklevel{4}\enspace 軸に平行な道（折れ線）の線積分を、ふつうの定積分に直して計算できる。
\item \bklevel{5}\enspace 場合分けの表を使って、ななめ・曲線・円の道、保存力、閉じた道の線積分を計算できる。
\end{itemize}
\end{goal}

\begin{hatena}
\noindent\begin{minipage}[c]{0.53\linewidth}
\raggedright
場所 $(x,\ y)$（m）ごとに向きと大きさのちがう力 $\vec F=(-y,\ x)$（N）がはたらく平面で、物を原点から点 $(1,\ 1)$ まで動かします。道Aは右へ1\,m、上へ1\,m。道Bは上へ1\,m、右へ1\,m。ななめの道もあります。\par\vspace{4pt}
\textbf{この力がする仕事は、道によって変わるでしょうか。遠回りすると、仕事は変わる？}
\end{minipage}\hfill
\begin{minipage}[c]{0.44\linewidth}
\centering
\begin{tikzpicture}[x=0.75mm, y=0.75mm]
  \draw[bkLine, line width=0.4pt, step=20] (0,0) grid (40,40);
  \kviiihfield{-\kviiihj}{\kviiihi}{40}{9}{1}
  \draw[bkOrange!85!black, line width=1.6pt, -{Stealth[length=2.4mm]}] (0,0) -- (40,0) -- (40,40);
  \draw[bkGreen!60!black, line width=1.6pt, dash pattern=on 3pt off 1.5pt, -{Stealth[length=2.4mm]}] (0,0) -- (0,40) -- (40,40);
  \draw[bkPurple!80!black, line width=1.3pt, dash pattern=on 1pt off 1.2pt] (0,0) -- (40,40);
  \node[font=\footnotesize, anchor=north] at (20,-1.5) {道A};
  \node[font=\footnotesize, anchor=east] at (-1.5,20) {道B};
  \node[font=\footnotesize, anchor=north] at (40,-1.5) {$(1,\ 0)$};
  \node[font=\footnotesize, anchor=south] at (0,41.5) {$(0,\ 1)$};
  \node[font=\footnotesize, anchor=west] at (44,38) {$(1,\ 1)$};
\end{tikzpicture}\par\vspace{2pt}
{\footnotesize 灰色の矢印が力 $\vec F$。オレンジが道A、\\緑の点線が道B、紫の点線がななめ。}
\end{minipage}
\end{hatena}

\noindent 4-12 で、仕事は「力と、動いた矢印の内積」だと学びました。5-12 では、力が場所で変わるときは、細かく分けて、かけて、たしました。でもそれは、まっすぐな道の上の話でした。

この章では、力の向きと大きさが場所ごとにちがい（8-5 のベクトル場）、道も曲がっているときの仕事を求めます。道を小さな1歩に分け、1歩ずつ「力と1歩の内積」をたします。この操作の名前が\textbf{線積分}です。8-6 と同じく、名前を、ふつうの定積分に直して計算します。どう直すかは道の形で決まるので、場合ごとに見ていきます。

この章の計算は、物理の本（力-4）の「仕事とエネルギー」、電気の「電圧」「起電力」、熱の「気体がする仕事」に、そのまま使われます。

%% ---------------- やってみよう・思い出そう ----------------
\begin{yattemiyou}[1歩ずつ「力×進んだ距離」をたそう]
道Aを、0.25\,m ずつの8歩に分けます。1歩ごとに、その1歩のまん中の場所で力 $\vec F=(-y,\ x)$ を計算し、1歩の移動 $(\Delta x,\ \Delta y)$ との内積 $F_x\Delta x+F_y\Delta y$ を求めて、たします。
\begin{enumerate}[label={\bkcircled[bkGreen]{\arabic*}}, leftmargin=2.4\zw, itemsep=1pt, topsep=3pt]
\item A列（A2〜A9）に、歩の番号1〜8を入れる。
\item B列・C列に、その1歩のまん中の $x$・$y$ を入れる。右へ進む4歩は $y=0$、上へ進む4歩は $x=1$。
\item D列・E列に、1歩の移動 $\Delta x$・$\Delta y$ を入れる。右へ進む歩は $(0.25,\ 0)$、上へ進む歩は $(0,\ 0.25)$。
\item F2 に \texttt{=-C2*D2+B2*E2} と入れて、F9 までコピーする。
\item F10 に \texttt{=SUM(F2:F9)} と入れる。道B・ななめの道でも同じ表を作ろう。
\end{enumerate}
\end{yattemiyou}

\noindent ④の式は、$F_x\Delta x+F_y\Delta y$ に $F_x=-y$（C列）、$F_y=x$（B列）を入れたものです（4-12 の、成分で計算する内積）。力は場所で変わるので、1歩の中のまん中の場所の力を使います（8-6 のタイルのまん中と同じ考え）。

\begin{bkcenter}
\bksheet{6}{15}{%
1 & 歩 & まん中の $x$ & まん中の $y$ & $\Delta x$ & $\Delta y$ & 内積 \\
2 & 1 & 0.125 & 0 & 0.25 & 0 & 0 \\
3 & 2 & 0.375 & 0 & 0.25 & 0 & 0 \\
4 & 3 & 0.625 & 0 & 0.25 & 0 & 0 \\
5 & 4 & 0.875 & 0 & 0.25 & 0 & 0 \\
6 & 5 & 1 & 0.125 & 0 & 0.25 & 0.25 \\
7 & 6 & 1 & 0.375 & 0 & 0.25 & 0.25 \\
8 & 7 & 1 & 0.625 & 0 & 0.25 & 0.25 \\
9 & 8 & 1 & 0.875 & 0 & 0.25 & 0.25 \\
10 & & & & & 合計 & 1 \\}
\bktabcaption{道Aの表計算。右へ進む間は $y=0$ なので $F_x=0$ で、仕事は0。上へ進む間は $F_y=x=1$ で、1歩ごとに $1\times0.25=0.25$ J。}
\label{bk:8-7:tab-sheet}
\end{bkcenter}

\begin{omoidasou}
\bkrecall{4-12}{仕事は、力 $\vec F$ と動いた矢印 $\vec s$ の内積 $\vec F\cdot\vec s=F_xs_x+F_ys_y$。力と動きが直角なら0。}
\bkrecall{5-12}{力が場所で変わるときの仕事は、細かく分けて「力×進んだ距離」をたした先 $\int F\,\dd x$。}
\bkrecall{4-11}{点 $\vec p_0$ から、向き $\vec v$ にまっすぐ進む点は $\vec p=\vec p_0+t\vec v$。}
\bkrecall{5-8}{短い時間 $\Delta t$ の間の位置の変化は、ほぼ「速度×$\Delta t$」。}
\bkrecall{8-3}{道 $x(t)$、$y(t)$ を歩くと、$\dfrac{\dd f}{\dd t}=\dfrac{\partial f}{\partial x}\dfrac{\dd x}{\dd t}+\dfrac{\partial f}{\partial y}\dfrac{\dd y}{\dd t}$（連鎖律）。}
\bkrecall{8-6}{積分の記号は操作の名前。名前の段→直す段→楽をする段で、ふつうの定積分に直して計算する。}
\end{omoidasou}



%% ---------------- 数学の道具 ----------------
\bksection{doguu}{数学の道具}

\bkstage{me}{目でわかる}{道を小さな1歩に分ける}{\bklevel{3}}

\bksubsection{1}{1歩ずつ、力と1歩の内積をたす}

\label{bk:8-7:ippo}物が道にそって動くとき、道を小さな移動 $\Delta\vec r_i$ に分けます。1つの区間の間は、力はほぼ一定の $\vec F_i$ とみなせます。その区間の仕事は、内積 $\vec F_i\cdot\Delta\vec r_i$ です（4-12）。全部の区間をたすと、道全体の仕事が近い値で出ます。

\begin{bkfigure}
\begin{tikzpicture}[x=1mm, y=1mm]
  \draw[bkInk, line width=1pt] plot[smooth] coordinates {(0,0) (20,14) (45,20) (70,16) (95,28)};
  \foreach \kviiihp/\kviiihq in {{(0,0)}/{(20,14)},{(20,14)}/{(45,20)},{(45,20)}/{(70,16)},{(70,16)}/{(95,28)}} {\draw[bkTheme, line width=1.4pt, -{Stealth[length=2.4mm]}] \kviiihp -- \kviiihq;}
  \foreach \kviiihp/\kviiihd in {{(10,7)}/{(12,-6)},{(32.5,17)}/{(10,10)},{(57.5,18)}/{(-4,13)},{(82.5,22)}/{(-12,8)}} {\draw[bkOrange!85!black, line width=1.3pt, -{Stealth[length=2.2mm]}] \kviiihp -- ++\kviiihd;}
  \fill[bkInk] (0,0) circle (0.9) (95,28) circle (0.9);
  \node[font=\footnotesize, anchor=north] at (0,-1.5) {はじめ};
  \node[font=\footnotesize, anchor=south] at (95,29.5) {おわり};
  \node[font=\footnotesize, anchor=north west, text=bkInk] at (45,12) {紺：1歩の移動 $\Delta\vec r_i$};
  \node[font=\footnotesize, anchor=north west, text=bkInk] at (45,6) {オレンジ：その場所の力 $\vec F_i$};
\end{tikzpicture}
\bkfigcaption{道を4つの小さな移動（紺）に分けた。1歩ごとの仕事は、その場所の力（オレンジ）と移動の内積。力が移動と同じ向きなら正、反対向きなら負、直角なら0。}
\label{bk:8-7:fig-ippo}
\end{bkfigure}

やってみようの表計算は、この考えを道Aで行ったものでした。道Bとななめの道でも同じように計算すると、次の表のようになります。くらべるために、もう1つの力 $\vec F_1=(y,\ x)$ でも計算しました（はてな？の力を $\vec F_2=(-y,\ x)$ とよびます。物理の本の力-4 と同じ名前です）。

\begin{bkcenter}
\begin{bktable}{colspec={Q[l,m,40mm] Q[c,m,26mm] Q[c,m,26mm] Q[c,m,26mm]}, row{2-Z}={font=\small}}
力 & 道A & 道B & ななめ \\
$\vec F_2=(-y,\ x)$（はてな？の力） & 1 J & $-1$ J & 0 J \\
$\vec F_1=(y,\ x)$ & 1 J & 1 J & 1 J \\
\end{bktable}
\bktabcaption{3つの道の仕事（0.25\,m ずつの表計算。どれもまん中の値でぴったり出る）。}
\label{bk:8-7:tab-michi}
\end{bkcenter}

道Bでは、はじめに上へ進む間は $x=0$ なので $F_y=0$ で、仕事は0です。そのあと右へ進む間は $y=1$ なので $F_x=-1$ で、力は進む向きと反対です。だから仕事は $-1\times1=-1$ J です。

\kviiihhead{気づくこと}
\begin{itemize}[leftmargin=1.2\zw, itemsep=2pt, topsep=0pt, label={\color{\bkTT}●}]
\item はてな？の力 $\vec F_2$ では、同じ点から同じ点へ動かしても、\textbf{道によって仕事が変わる}（1、$-1$、0）。
\item $\vec F_1$ では、どの道でも仕事は1。\textbf{道によらない}力もある。
\end{itemize}
どんな力なら道によらないのかは、⑨で調べます。

\bksubsection{2}{名前の段：線積分の記号を読む}

\label{bk:8-7:namae}1歩を細かくした先の値を、次のように書きます。物理の本（力-4）と同じ書き方です。
\kviiihframe[104mm]{\Large $\displaystyle\int_C\vec F\cdot\dd\vec r=\lim\sum_i\vec F_i\cdot\Delta\vec r_i$\\[3pt]\normalsize\sffamily 「道 $C$ にそって、力と1歩の矢印の内積をたし合わせた先」}
これを、道 $C$ にそった\textbf{線積分}（せんせきぶん）といいます。\bkidx{せんせきぶん}{線積分}$\lim$ は、区切りを細かくした先、という意味です。

\kviiihdanhead{1}{記号を読む：何を・どの道で・1歩の量}
\begin{enumerate}[label={\bkcircled[\bkTT]{\arabic*}}, leftmargin=2.4\zw, itemsep=2pt, topsep=2pt]
\item \textbf{どの道で}：∫ の右下の $C$ は、物が通った道です（英語の curve（曲線）の頭文字）。道には\textbf{向き}があります。はじめの点からおわりの点へ進みます。
\item \textbf{何を}：$\vec F$ は、場所ごとにちがう力（ベクトル場、8-5）です。
\item \textbf{1歩の量}：$\dd\vec r$ は、とても小さな1歩の移動の矢印です。$\vec F\cdot\dd\vec r$ は、その1歩の仕事（内積）です。
\end{enumerate}
道の上の場所は、「はじめから何 m 進んだ所か」や「何秒後の場所か」のように、1つの数で決まります。だから ∫ は1つです（8-6 の「∫ の数＝場所を表す数の数」）。

\kviiihhead{記号は「操作の名前」}
8-6 の重積分と同じく、$\int_C\vec F\cdot\dd\vec r$ は操作の\textbf{名前}です。$C$ という名前と、$\dd\vec r$ という矢印しか書いてないので、このままでは計算できません。ふつうの定積分に\textbf{直す}と、次のような形になります（どう直すかは④・⑥で）。
\[ \kviiihNt{\int_C\vec F\cdot\dd\vec r}=\kviiihSt{\int_0^1F_x(x,\ 0)\,\dd x+\int_0^1F_y(1,\ y)\,\dd y} \]
これは道Aの場合です。道が変われば、直した形も変わります。

\begin{taisetsu}
\begin{itemize}[leftmargin=1.2\zw, itemsep=2pt, topsep=0pt, label={\color{\bkTT}●}]
\item 線積分 $\int_C\vec F\cdot\dd\vec r$ は、「道 $C$ を小さな1歩に分け、力と1歩の内積をたし、細かくした先」という操作の名前。
\item 道には向きがある。力によっては、同じ2点を結んでも、道によって値が変わる。
\end{itemize}
\end{taisetsu}

\bkstage{te}{手でわかる}{表計算と、軸に平行な道}{\bklevel{4}}

\bksubsection{3}{表計算で、細かくしてたす}

\label{bk:8-7:hyou}道が曲がっているときも、表計算で近い値が出せます。半径1\,m の円を、左回りに1周する道で、はてな？の力 $\vec F_2$ の仕事を求めます。円を $n$ 本の短い線分（弦）に分け、それぞれのまん中の点の力と、線分の矢印の内積をたします（物理の本の力-4 の Python と同じ方法）。

\begin{bkfigure}
\begin{tikzpicture}[x=1mm, y=1mm]
  \foreach \kviiihn/\kviiihx in {4/0,8/42,12/84} {%
    \begin{scope}[xshift=\kviiihx mm]
      \draw[bkGray, line width=0.5pt, dash pattern=on 1.2pt off 1pt] (16,16) circle (14);
      \draw[bkTheme, line width=1.1pt] plot[samples=\kviiihn+1, domain=0:360] ({16+14*cos(\x)},{16+14*sin(\x)});
      \fill[bkInk] (16,16) circle (0.7);
      \node[font=\footnotesize, anchor=north] at (16,0) {$n=\kviiihn$};
    \end{scope}}
\end{tikzpicture}
\bkfigcaption{円（点線）を、$n$ 本の線分（紺）に分ける。$n$ を大きくすると、線分の道は円の道に近づく。}
\label{bk:8-7:fig-takakkei}
\end{bkfigure}

\begin{bkcenter}
\begin{bktable}{colspec={Q[c,m,30mm] Q[c,m,22mm] Q[c,m,22mm] Q[c,m,22mm] Q[c,m,22mm]}, row{2-Z}={font=\small}}
線分の数 $n$ & 4 & 8 & 12 & 100 \\
仕事の合計（J） & 4 & 5.657 & 6.000 & 6.279 \\
\end{bktable}
\bktabcaption{$n$ をふやすと、$2\pi\approx6.283$ に近づく。物理の本の Python の答え 6.2791 は $n=100$ のときの値。}
\label{bk:8-7:tab-n}
\end{bkcenter}
$n$ を大きくした先は $2\pi$ になります。このことは、⑥で、ふつうの定積分に直して確かめます。

\kviiihhead{Python でたす}
同じ計算を Python で行うと、次のようになります（物理の本の力-4 と同じ方法）。円の上に $n+1$ 個の点をとり、となりどうしの点を結ぶ線分のまん中の点と、線分の矢印を、配列でまとめて計算します。
\Needspace{11\baselineskip}
\begin{bkcode}[Python]
import numpy as np
n = 100
th = np.linspace(0, 2*np.pi, n+1)   \# 円の上の n+1 個の点の角
x, y = np.cos(th), np.sin(th)
mx = (x[:-1] + x[1:]) / 2           \# 線分のまん中の点
my = (y[:-1] + y[1:]) / 2
dx, dy = np.diff(x), np.diff(y)     \# 1歩の矢印（Δx, Δy）
W = np.sum(-my*dx + mx*dy)          \# F2=(-y, x) との内積をたす
print(round(W, 4))
\end{bkcode}
\begin{bkcodeout}[実行結果]
6.2791
\end{bkcodeout}
\texttt{n = 1000} にすると 6.2832 になり、$2\pi=6.28318\dots$ にさらに近づきます。

\bksubsection{4}{直す段：軸に平行な道}

\label{bk:8-7:jiku}道Aや道Bのように、道が「横だけ」「たてだけ」の線分でできているときは、ふつうの定積分にすぐ直せます。

\kviiihdanhead{2}{動く文字は1つだけ}
右へ進む区間では、$y$ は変わりません（$\dd y=0$）。だから1歩の仕事は次のとおりです。
\kviiihline{$\vec F\cdot\dd\vec r=F_x\,\dd x+F_y\times0=F_x\,\dd x$}
上へ進む区間では、$x$ は変わらず（$\dd x=0$）、$F_y\,\dd y$ です。
\begin{enumerate}[label={\bkcircled[\bkTT]{\arabic*}}, leftmargin=2.4\zw, itemsep=2pt, topsep=2pt]
\item 道を、横の区間とたての区間に分ける。
\item 横の区間：$y$ をその区間の値で止めて、$F_x$ を $x$ で積分する。
\item たての区間：$x$ を止めて、$F_y$ を $y$ で積分する。
\item 範囲は「はじめの値から、おわりの値まで」。左や下へ進む区間では、下のはしの数が上のはしより大きくなる。
\item 区間ごとの答えをたす。
\end{enumerate}

\kviiihhead{道A（$\vec F_2=(-y,\ x)$）}
\begin{bkhenkei}[0.56]
\int_0^1F_x(x,\ 0)\,\dd x=\int_0^1(-0)\,\dd x=0 & 右へ：$y=0$ に止める。$F_x=-y=0$ \\
\int_0^1F_y(1,\ y)\,\dd y=\int_0^1 1\,\dd y=1 & 上へ：$x=1$ に止める。$F_y=x=1$ \\
0+1=1 & 2つの区間をたす \\
\end{bkhenkei}

\kviiihhead{道B（$\vec F_2=(-y,\ x)$）}
\begin{bkhenkei}[0.56]
\int_0^1F_y(0,\ y)\,\dd y=\int_0^1 0\,\dd y=0 & 上へ：$x=0$ に止める。$F_y=x=0$ \\
\int_0^1F_x(x,\ 1)\,\dd x=\int_0^1(-1)\,\dd x=-1 & 右へ：$y=1$ に止める。$F_x=-y=-1$ \\
0+(-1)=-1 & 2つの区間をたす \\
\end{bkhenkei}
表計算と同じ1と $-1$ が出ました。

\kviiihhead{逆向きに進むと、符号が変わる}
道Aを逆に、$(1,\ 1)$ から $(1,\ 0)$ へ下り、$(0,\ 0)$ へ左に進むとします。下る区間は $\int_1^0F_y(1,\ y)\,\dd y=\int_1^0 1\,\dd y=-1$、左へ進む区間は $\int_1^0 0\,\dd x=0$ で、合計は $-1$ です。同じ道を逆向きに進むと、1歩の矢印 $\dd\vec r$ がすべて反対向きになるので、線積分の符号が変わります。だから、道の\textbf{向き}は、いつも書いておきます。

\bkkokomade[道を小さな1歩に分けて、力と1歩の内積を表計算でたし、軸に平行な道なら、ふつうの定積分に直して計算できれば十分です。この先は、ななめ・曲線・円の道と、計算を楽にする方法です。]

\Needspace{9\baselineskip}\bkstage{shiki}{式でわかる}{道の形ごとに、直し方を選ぶ}{\bklevel{5}}

\bksubsection{5}{場合分けの表と流れ図}

\label{bk:8-7:basho}道がななめや曲線になると、$x$ と $y$ がいっしょに変わります。そのときは、道の上の場所を\textbf{1つの数}（たとえば時刻 $t$）で表して、ふつうの定積分1回に直します。どの数を使うかは、道の形で決まります。8-6 と同じく、「こんなときは、こう直す」を表にしました（記号 A〜K）。

\kviiihcasetitle{8-7}{線積分：道の形で、どう直す？}
\begin{bkcenter}
\begin{bktable}{colspec={Q[c,m,10mm] Q[l,m,30mm] X[l,m] Q[c,m,12mm]}, row{2-Z}={font=\small}, rowsep=3pt, colsep=4pt}
記号 & こんなとき & 手順（ふつうの定積分に直す） & 例題 \\
\kviiihcase{A} & 道が軸に平行な線分（横だけ・たてだけ）でできている & 区間ごとに、動く文字は1つだけ。横は $F_x\,\dd x$（$y$ は止める）、たては $F_y\,\dd y$（$x$ は止める） & 1・8 \\
\kviiihcase{B} & まっすぐなななめの道（点 A から点 B へ） & $\vec r(t)=\vec A+t(\vec B-\vec A)$、$0\leqq t\leqq1$。$\dd\vec r=(\vec B-\vec A)\,\dd t$ & 2 \\
\kviiihcase{C} & 道が $y=g(x)$ のグラフ & $x$ を道の数にする。$\dd\vec r=(1,\ g'(x))\,\dd x$。$y$ には $g(x)$ を入れる & 3 \\
\kviiihcase{D} & 道が円・円弧 & 角 $\theta$ を道の数にする。$\vec r=(R\cos\theta,\ R\sin\theta)$、$\dd\vec r=(-R\sin\theta,\ R\cos\theta)\,\dd\theta$ & 4・9 \\
\end{bktable}
\end{bkcenter}

\kviiihcasetitle{8-7}{（つづき）区間・長さ・楽をする場合}
\begin{bkcenter}
\begin{bktable}{colspec={Q[c,m,10mm] Q[l,m,30mm] X[l,m] Q[c,m,12mm]}, row{2-Z}={font=\small}, rowsep=3pt, colsep=4pt}
記号 & こんなとき & 手順 & 例題 \\
\kviiihcase{E} & ななめ・曲線のある折れ線 & 区間ごとに A〜D で計算して、たす。逆向きの区間は符号が変わる & 5 \\
\kviiihcase{I} & 長さでたす $\int_C f\,\dd s$（向きのない量） & $\dd s=|\vec r\,'(t)|\,\dd t$（三平方）。向きを逆にしても同じ値 & 9 \\
\kviiihcase{F} & 力が式で書いてなく、大きさと向きで決まる（いつも道と直角、または同じ向き・反対向きで大きさ一定。摩擦力など） & 内積を先に：直角なら0。平行なら $\pm F\times$（道の長さ） & 6 \\
\kviiihcase{G} & $U$ のわかる保存力（$\vec F=-\mathrm{grad}\,U$） & 道によらない：$U(\text{はじめ})-U(\text{おわり})$。楽な道に取りかえてもよい & 7 \\
\kviiihcase{H} & 閉じた道（∮） & 保存力なら0。そうでなければ1周を F〜D の問いで計算（8-11 では「rot×面積」） & 4・5・8 \\
\kviiihcase{J} & データだけ、または数値で近い値を求める & 表計算で $\sum\vec F_i\cdot\Delta\vec r_i$（まん中の点で） & 10 \\
\kviiihcase{K} & 状態の図（$p$-$V$ 図）の上の道 & $\int p\,\dd V$ も道にそってたす。区間ごとに A & 11 \\
\end{bktable}
\end{bkcenter}

A〜E と I が②の\textbf{直す段}、F・G・H・J・K が③の\textbf{楽をする段}です（K は物理で使う形で、区間ごとの直し方は A と同じ）。どの順で問えばよいかを、2つの流れ図にしました。先に楽をする段の問い（図\ref{bk:8-7:fig-flow}）、そのあとに直す段の道の形の問い（図\ref{bk:8-7:fig-flow2}）をたどります。長さでたす積分（I）でも、道の形の問い（A〜E）は同じです。

\begin{bkfigure}
\begin{tikzpicture}[x=1mm, y=1mm]
  % 流れ図の組み方（Ver2 の共通の決まり 2章）：問いは左の列にたて1列。「はい」は右へ、「いいえ」は下へ。線は交わらない
  \node[kviiih do, anchor=north west] (s) at (4,0) {記号を読む（名前の段）：何を・どの道で（向き）\\長さでたす $f\,\dd s$ なら $\dd s=|\vec r\,'|\,\dd t$ にする\enspace\kviiihcase{I}};
  \node[kviiih q, aspect=3.6, anchor=north] (q1) at ([yshift=-8mm]s.south -| 29,0) {データだけ・数値で？};
  \node[kviiih do, anchor=west, minimum width=60mm] (rJ) at (q1 -| 72,0) {表計算で $\sum\vec F_i\cdot\Delta\vec r_i$\enspace\kviiihcase{J}};
  \node[kviiih q, aspect=3.6, anchor=north] (qK) at ([yshift=-8mm]q1.south) {$p$-$V$ 図の上の道？};
  \node[kviiih do, anchor=west, minimum width=60mm] (rK) at (qK -| 72,0) {区間ごとに $\int p\,\dd V$\enspace\kviiihcase{K}};
  \node[kviiih q, aspect=3.6, anchor=north] (q2) at ([yshift=-8mm]qK.south) {$U$ のわかる保存力？};
  \node[kviiih do, anchor=west, minimum width=60mm] (rG) at (q2 -| 72,0) {$U(\text{はじめ})-U(\text{おわり})$\enspace\kviiihcase{G}};
  \node[kviiih q, aspect=3.6, anchor=north] (q3) at ([yshift=-8mm]q2.south) {閉じた道（∮）？};
  \node[kviiih do, anchor=west, minimum width=60mm] (rH) at (q3 -| 72,0) {1周を下の問いで計算\\（保存力なら0）\enspace\kviiihcase{H}};
  \node[kviiih q, aspect=3.6, anchor=north] (q4) at ([yshift=-17mm]q3.south) {力が式で書いてない？};
  \node[kviiih do, anchor=west, minimum width=60mm] (rF) at (q4 -| 72,0) {0 か $\pm F\times$長さ\enspace\kviiihcase{F}};
  \node[kviiih do, anchor=north west, minimum width=128mm] (x) at ([yshift=-8mm]q4.south -| 4,0) {直す段へ（図\ref{bk:8-7:fig-flow2}）：道の形の問いで、ふつうの定積分に直す};
  \node[kviiih note, anchor=north west] at ([yshift=-3mm]x.south west) {\textbf{確かめ}（答えが出たら）：逆向きで符号が変わるか。保存力なら、別の道で同じか};
  \draw[kviiih go] (s.south -| q1) -- (q1.north);
  \foreach \kviiihq/\kviiihr in {q1/rJ,qK/rK,q2/rG,q3/rH,q4/rF} {\draw[kviiih go] (\kviiihq.east) -- node[kviiih yes] {はい} (\kviiihq.east -| \kviiihr.west);}
  \foreach \kviiihq/\kviiihr in {q1/qK,qK/q2,q2/q3} {\draw[kviiih go] (\kviiihq.south) -- node[kviiih no] {いいえ} (\kviiihr.north);}
  \draw[kviiih go] (q3.south) -- node[kviiih no, pos=0.3] {いいえ} (q4.north);
  \draw[kviiih go] (q4.south) -- node[kviiih no] {いいえ} (q4.south |- x.north);
  % H：1周を、次の問い（力が式で書いてない？）からたどる
  \draw[kviiih go, shorten >=0pt] (rH.south) |- ($(q3.south)+(0.9mm,-11mm)$);
\end{tikzpicture}
\bkfigcaption{線積分の流れ図（1）楽をする段。ひし形は問い、四角はすること。「はい」は右へ、「いいえ」は下へ進む。H の四角からは、矢印の先の問いへ。}
\label{bk:8-7:fig-flow}
\end{bkfigure}

\begin{bkfigure}
\begin{tikzpicture}[x=1mm, y=1mm]
  \node[kviiih do, anchor=north] (s2) at (29,0) {図\ref{bk:8-7:fig-flow}の「直す段へ」から};
  \node[kviiih q, aspect=2.6, anchor=north] (q5) at ([yshift=-8mm]s2.south) {ななめ・曲線の\\ある折れ線？};
  \node[kviiih do, anchor=west, minimum width=60mm] (rE) at (q5 -| 72,0) {区間ごとに下の問いで\\計算して、たす\enspace\kviiihcase{E}};
  \node[kviiih q, aspect=3.6, anchor=north] (q6) at ([yshift=-17mm]q5.south) {軸に平行な線分だけ？};
  \node[kviiih do, anchor=west, minimum width=60mm] (rA) at (q6 -| 72,0) {$F_x\,\dd x$ か $F_y\,\dd y$\enspace\kviiihcase{A}};
  \node[kviiih q, aspect=3.6, anchor=north] (q7) at ([yshift=-8mm]q6.south) {まっすぐなななめ？};
  \node[kviiih do, anchor=west, minimum width=60mm] (rB) at (q7 -| 72,0) {$\vec A+t(\vec B-\vec A)$\enspace\kviiihcase{B}};
  \node[kviiih q, aspect=3.6, anchor=north] (q8) at ([yshift=-8mm]q7.south) {$y=g(x)$ のグラフ？};
  \node[kviiih do, anchor=west, minimum width=60mm] (rC) at (q8 -| 72,0) {$(1,\ g')\,\dd x$\enspace\kviiihcase{C}};
  \node[kviiih q, aspect=3.6, anchor=north] (q9) at ([yshift=-8mm]q8.south) {円・円弧？};
  \node[kviiih do, anchor=west, minimum width=60mm] (rD) at (q9 -| 72,0) {角 $\theta$ で\enspace\kviiihcase{D}};
  \node[kviiih do, anchor=north] (rT) at ([yshift=-8mm]q9.south) {道を $\vec r(t)$ で表し、$\dd\vec r=\vec r\,'(t)\,\dd t$};
  \node[kviiih note, anchor=north west] at ([yshift=-4mm]rT.south -| 4,0) {\textbf{確かめ}（答えが出たら）：逆向きで符号が変わるか。保存力なら、別の道で同じか};
  \draw[kviiih go] (s2.south) -- (q5.north);
  \foreach \kviiihq/\kviiihr in {q5/rE,q6/rA,q7/rB,q8/rC,q9/rD} {\draw[kviiih go] (\kviiihq.east) -- node[kviiih yes] {はい} (\kviiihq.east -| \kviiihr.west);}
  \draw[kviiih go] (q5.south) -- node[kviiih no, pos=0.3] {いいえ} (q6.north);
  \foreach \kviiihq/\kviiihr in {q6/q7,q7/q8,q8/q9,q9/rT} {\draw[kviiih go] (\kviiihq.south) -- node[kviiih no] {いいえ} (\kviiihr.north);}
  % E：区間ごとに、次の問い（軸に平行な線分だけ？）からたどる
  \draw[kviiih go, shorten >=0pt] (rE.south) |- ($(q5.south)+(0.9mm,-11mm)$);
\end{tikzpicture}
\bkfigcaption{線積分の流れ図（2）直す段。図\ref{bk:8-7:fig-flow}の「直す段へ」から、道の形の問いで直し方を選ぶ。E の四角からは、矢印の先の問いへ。}
\label{bk:8-7:fig-flow2}
\end{bkfigure}

\kviiihhead{流れ図の使い方：道B}
はてな？の力 $\vec F_2$ で、道Bの線積分をたどります。
\begin{enumerate}[label={\bkcircled[\bkTT]{\arabic*}}, leftmargin=2.4\zw, itemsep=1pt, topsep=2pt]
\item 名前の段：何を＝$\vec F_2\cdot\dd\vec r$、道＝道B（$(0,\ 0)\to(0,\ 1)\to(1,\ 1)$）。
\item データだけ・数値で？\enspace いいえ。\enspace $p$-$V$ 図の上の道？\enspace いいえ。\enspace $U$ のわかる保存力？\enspace いいえ（$\vec F_2$ は道で値が変わった。表\ref{bk:8-7:tab-michi}）。
\item 閉じた道？\enspace いいえ。\enspace 力が式で書いてない？\enspace いいえ（式 $(-y,\ x)$ で書いてある）。\enspace 「直す段へ」から図\ref{bk:8-7:fig-flow2}へ。
\item ななめ・曲線のある折れ線？\enspace いいえ。\enspace 軸に平行な線分だけ？\enspace はい。
\end{enumerate}
こうして場合 \kviiihcase{A} に決まります。④でしたとおりです。1回で全部の場合を覚えなくても大丈夫です。問題を解くたびに、表と流れ図にもどって確かめましょう。

\kviiihhead{どの場合かな？}
次の線積分は、場合分けの表8-7のどの場合でしょう。流れ図を上からたどって、先に自分で考えてから答えを見ましょう。
\begin{enumerate}[label={\bkcircled[\bkTT]{\arabic*}}, leftmargin=2.4\zw, itemsep=3pt, topsep=2pt]
\item 原点から $(2,\ 0)$ を通って $(2,\ 3)$ まで、$\vec F=(x,\ y^2)$\par\nobreak\noindent 答え：\kviiihcase{A}。軸に平行な線分だけ。
\item $(1,\ 1)$ から $(3,\ 2)$ へまっすぐ、$\vec F=(y,\ x^2)$\par\nobreak\noindent 答え：\kviiihcase{B}。まっすぐなななめ。
\item ばねの力（$U=\tfrac{1}{2}kx^2$）が、$x=0.1$ から $x=0.3$ までにする仕事\par\nobreak\noindent 答え：\kviiihcase{G}。$U$ のわかる保存力。
\item 机の上で物をおして1周させるときの、動摩擦力の仕事\par\nobreak\noindent 答え：\kviiihcase{H}→\kviiihcase{F}。閉じた道で、摩擦力は式で書いてなく、大きさと向きで決まる。
\item 道の上の10か所ではかった力のデータから仕事を出す\par\nobreak\noindent 答え：\kviiihcase{J}。データだけ。
\end{enumerate}

\bksubsection{6}{直す段：道を1つの数 $t$ で表す}

\label{bk:8-7:t}道がななめや曲線のときは、道の上の場所を、1つの数 $t$ で表します。$t$ は「道を歩く人の時刻」と思ってかまいません。時刻 $t$ に場所 $\vec r(t)=(x(t),\ y(t))$ にいて、$t=a$ ではじめの点、$t=b$ でおわりの点に着くとします。

\kviiihdanhead{2}{1歩の矢印は $\dd\vec r=\vec r\,'(t)\,\dd t$}
短い時間 $\dd t$ の間に動く矢印は、ほぼ「速度×時間」です（5-8 の近似）。速度は位置の矢印を微分したもの $\vec r\,'(t)=(x'(t),\ y'(t))$ なので、
\kviiihframe[110mm]{\Large $\dd\vec r=\vec r\,'(t)\,\dd t$}
\[ \kviiihNt{\int_C\vec F\cdot\dd\vec r}=\kviiihSt{\int_a^b\vec F(\vec r(t))\cdot\vec r\,'(t)\,\dd t} \]
右辺は、$t$ だけの関数を $a$ から $b$ まで積分する、ふつうの定積分です。手順は次の4つです。
\begin{enumerate}[label={\bkcircled[\bkTT]{\arabic*}}, leftmargin=2.4\zw, itemsep=2pt, topsep=2pt]
\item 道を $\vec r(t)$ で表し、はじめとおわりの $t$ を決める。
\item $\vec r\,'(t)$ を計算する（成分ごとに $t$ で微分）。
\item $\vec F$ の $x$・$y$ に、$x(t)$・$y(t)$ を入れる（$\vec F$ を $t$ の式にする）。
\item 内積 $\vec F\cdot\vec r\,'(t)$ を計算し、$t$ で $a$ から $b$ まで積分する。
\end{enumerate}

\begin{bkfigure}
\begin{tikzpicture}[x=1mm, y=1mm]
  \draw[bkInk, line width=1pt] plot[smooth, domain=0:1, samples=30] ({10+80*\x},{6+30*\x*\x});
  \foreach \kviiiht in {0,0.25,0.5,0.75,1} {\fill[bkInk] ({10+80*\kviiiht},{6+30*\kviiiht*\kviiiht}) circle (0.9);
    \node[font=\footnotesize, anchor=north west] at ({11+80*\kviiiht},{5.5+30*\kviiiht*\kviiiht}) {$t=\kviiiht$};}
  \draw[bkTheme, line width=1.4pt, -{Stealth[length=2.4mm]}] (50,13.5) -- ++(16,6);
  \node[font=\footnotesize, anchor=south east, text=bkInk, fill=white, inner sep=0.8pt] at (57,18) {$\vec r\,'(t)\,\dd t$};
\end{tikzpicture}
\bkfigcaption{道の上の場所に、時刻 $t$ の目もりをつける。時刻 $t$ の所から短い時間 $\dd t$ に動く矢印は、速度 $\vec r\,'(t)$ に $\dd t$ をかけたもの。}
\label{bk:8-7:fig-t}
\end{bkfigure}

\kviiihhead{場合 B：まっすぐなななめの道}
点 $\vec A$ から点 $\vec B$ へまっすぐ進む道は、4-11 の「まっすぐ一定の速さで進む点」と同じで、
\kviiihline{$\vec r(t)=\vec A+t(\vec B-\vec A)$\quad （$0\leqq t\leqq1$）、\qquad $\vec r\,'(t)=\vec B-\vec A$}
です。$t=0$ で $\vec A$、$t=1$ で $\vec B$ に着きます。はてな？のななめの道（$(0,\ 0)$ から $(1,\ 1)$）は $\vec r(t)=(t,\ t)$、$\vec r\,'(t)=(1,\ 1)$ です。$\vec F_2=(-y,\ x)$ に $x=t$、$y=t$ を入れると $(-t,\ t)$ で、内積は $-t\times1+t\times1=0$。だから仕事は0です（表\ref{bk:8-7:tab-michi}と同じ）。

\kviiihhead{場合 C：道が $y=g(x)$ のグラフ}
道がグラフ $y=g(x)$ のときは、$x$ そのものを道の数にすると楽です（$t$ のかわりに $x$）。
\kviiihline{$\vec r(x)=(x,\ g(x))$、\qquad $\dd\vec r=(1,\ g'(x))\,\dd x$}
大事なのは、$\vec F$ の中の $y$ に、$g(x)$ を入れることです。$y$ を残したまま $x$ で積分してはいけません（まちがえやすいところ③）。

\kviiihhead{場合 D：道が円・円弧}
半径 $R$ の円の上の場所は、角 $\theta$ を使って $\vec r(\theta)=(R\cos\theta,\ R\sin\theta)$ と書けます（4-15）。$\theta$ で微分すると、
\kviiihline{$\dd\vec r=(-R\sin\theta,\ R\cos\theta)\,\dd\theta$}
です。この矢印は、位置の矢印 $(R\cos\theta,\ R\sin\theta)$ を左に90° 回した向きで、円の接線の向きです（内積が $-R^2\cos\theta\sin\theta+R^2\sin\theta\cos\theta=0$ で、直角）。長さは $R\,\dd\theta$ で、4-2 の弧の長さです。

\kviiihhead{③の答えを確かめる：円を1周}
半径1の円を左回りに1周（$\theta$ が0から $2\pi$）する道で、$\vec F_2=(-y,\ x)$ の線積分を求めます。
\begin{bkhenkei}[0.52]
\vec F_2=(-\sin\theta,\ \cos\theta) & $x=\cos\theta$、$y=\sin\theta$ を入れる \\
\vec F_2\cdot\vec r\,'=\sin^2\theta+\cos^2\theta=1 & $\vec r\,'=(-\sin\theta,\ \cos\theta)$ との内積 \\
\int_0^{2\pi}1\,\dd\theta=2\pi & ふつうの定積分 \\
\end{bkhenkei}
表\ref{bk:8-7:tab-n}の値が近づいていった先は、$2\pi\approx6.283$ でした。$\vec F_2$ は、どこでも道と同じ向き（円の接線の向き）を向いていて、大きさは1です。だから、仕事は「1×円周の長さ $2\pi$」になります（これは⑨の場合 F の考え方でもあります）。

\bksubsection{7}{向きと、区間に分ける道}

\label{bk:8-7:muki}道 $C$ を逆向きに進む道を $-C$ と書くと、
\kviiihline{$\displaystyle\int_{-C}\vec F\cdot\dd\vec r=-\int_C\vec F\cdot\dd\vec r$}
です。1歩の矢印がすべて反対向きになるからです。$t$ で表した積分なら、はじめとおわりの $t$ を入れかえることになり、ふつうの定積分の性質 $\int_b^a=-\int_a^b$（5-11）からも同じことがわかります。

\kviiihhead{場合 E：折れ線は、区間ごとに直してたす}
道がいくつかの区間（線分・曲線）をつないだものなら、区間ごとに A〜D で直して、たします。重積分の場合 F（形の和）と同じ考えです。区間ごとに向きを確かめます。たとえば、道Aで $(1,\ 1)$ まで行き、ななめの道を通って原点へもどる三角形の道では、もどる区間は「ななめの道を逆向き」なので、ななめの道の値の符号を変えたものになります（例題5）。

\bksubsection{8}{長さでたす線積分：場合 I}

\label{bk:8-7:ds}道にそってたす量が、力のような矢印ではなく、ただの数のこともあります。たとえば、太さがちがう針金の質量は、「1\,m あたりの質量 $\lambda$×1歩の長さ」をたしたものです。これを
\kviiihline{$\displaystyle\int_C\lambda\,\dd s$\qquad （$\dd s$：1歩の長さ）}
と書きます（$s$ は道にそった長さの文字。4-2 の弧の長さ $s=r\theta$ と同じ）。1歩の長さは、1歩の矢印 $\dd\vec r=(x'\,\dd t,\ y'\,\dd t)$ の長さなので、三平方の定理（4-4）から
\kviiihframe[100mm]{\Large $\dd s=|\vec r\,'(t)|\,\dd t=\sqrt{x'(t)^2+y'(t)^2}\,\dd t$}
です。円なら $|\vec r\,'(\theta)|=\sqrt{R^2\sin^2\theta+R^2\cos^2\theta}=R$ なので、$\dd s=R\,\dd\theta$ です。

長さでたす線積分は、道を逆向きに進んでも値は変わりません。1歩の長さは、向きを変えても正のままだからです。矢印でたす $\int_C\vec F\cdot\dd\vec r$ との大きなちがいです。$\lambda=1$ なら、$\int_C\dd s$ は道の長さそのものです。

\Needspace{12\baselineskip}\bksubsection{9}{楽をする段：直角・平行、保存力、閉じた道}

\label{bk:8-7:raku}\kviiihdanhead{3}{計算する前に、楽な道をさがす}

\kviiihhead{場合 F：力が道と直角・平行}
\begin{itemize}[leftmargin=1.2\zw, itemsep=2pt, topsep=0pt, label={\color{\bkTT}●}]
\item 力が道といつも直角なら、どの1歩も内積が0なので、仕事は0です。円運動をする物を中心へ引く力（向心力）は、いつも道と直角なので、仕事をしません。
\item 力がいつも道と同じ向きで大きさ $F$ が一定なら、仕事は $F\times$（道の長さ）。反対向きなら $-F\times$（道の長さ）。床の上をすべらせる物にはたらく動摩擦力は、大きさ $\mu'N$ で、いつも動く向きと反対なので、仕事は $-\mu'N\times$（道の長さ）です。遠回りするほど、摩擦の仕事は大きく負になります（物理 力-4）。
\end{itemize}

\kviiihhead{場合 G：保存力なら、はじめとおわりだけで決まる}
8-4 で、位置エネルギー $U$ から力が $\vec F=-\mathrm{grad}\,U$ と作れる場合を学びました。このような力を\textbf{保存力}（ほぞんりょく）といいます。\bkidx{ほぞんりょく}{保存力}保存力の線積分は、道によらず、はじめとおわりの $U$ だけで決まります。
\kviiihframe[110mm]{\Large $\displaystyle\int_C\vec F\cdot\dd\vec r=U(\text{はじめ})-U(\text{おわり})$\\[2pt]\normalsize\sffamily （$\vec F=-\mathrm{grad}\,U$ のとき）}
\textbf{なぜか}：道を $\vec r(t)$（$a\leqq t\leqq b$）で表し、歩く人が感じる $U$ の変わる速さを、8-3 の連鎖律で計算します。
\begin{bkhenkei}[0.56]
\frac{\dd}{\dd t}U(x(t),\ y(t))=\frac{\partial U}{\partial x}x'+\frac{\partial U}{\partial y}y' & 8-3 の連鎖律 \\
=\mathrm{grad}\,U\cdot\vec r\,'(t)=-\vec F\cdot\vec r\,'(t) & 内積の形。$\mathrm{grad}\,U=-\vec F$ \\
\int_a^b\vec F\cdot\vec r\,'\,\dd t=-\Bigl[\,U(\vec r(t))\,\Bigr]_a^b & $-U$ を微分すると $\vec F\cdot\vec r\,'$（5-11） \\
=U(\text{はじめ})-U(\text{おわり}) & $t=a$ ではじめ、$t=b$ でおわり \\
\end{bkhenkei}
$\vec F_1=(y,\ x)$ は、$U=-xy$ とすると $-\mathrm{grad}\,U=-(-y,\ -x)=(y,\ x)$ なので、保存力です。$(0,\ 0)$ から $(1,\ 1)$ への仕事は、次のとおりです。
\kviiihline{$U(0,\ 0)-U(1,\ 1)=0-(-1)=1$}表\ref{bk:8-7:tab-michi}で、どの道でも1になったのは、このためです。重力も保存力で（$U=mgz$）、主役の山の $(30,\ 20)$（83\,m）から頂上（100\,m）まで、50\,kg の人がゆっくり登るときの仕事は、どの道でも 8330\,J です（例題7）。

\begin{bkfigure}
\begin{tikzpicture}[x=1mm, y=1mm]
  \foreach \kviiihh in {30,40,50,60,70,80,90} {\draw[bkGreen!60!black, line width=0.5pt] (0,0) ellipse ({0.5*sqrt((100-\kviiihh)/0.01)} and {0.5*sqrt((100-\kviiihh)/0.02)});}
  \fill[bkInk] (0,0) circle (1);
  \node[font=\footnotesize, anchor=north east, fill=white, inner sep=0.6pt] at (-1,-1.2) {頂上 100\,m};
  \fill[bkOrange!85!black] (15,10) circle (1);
  \node[font=\footnotesize, anchor=west, fill=white, inner sep=0.8pt] at (46,14) {$(30,\ 20)$：83\,m};
  \draw[bkInk, line width=0.4pt] (45.5,14) -- (16,10.4);
  \draw[bkOrange!85!black, line width=1.4pt, -{Stealth[length=2.4mm]}] (15,10) -- (1.3,0.9);
  \draw[bkBlue, line width=1.4pt, dash pattern=on 3pt off 1.5pt, -{Stealth[length=2.4mm]}] plot[smooth] coordinates {(15,10) (26,-6) (8,-24) (-18,-14) (-1.3,-0.9)};
  \node[font=\footnotesize, anchor=west, fill=white, inner sep=0.8pt] at (46,-12) {青の点線：遠回りの道};
  \node[font=\footnotesize, anchor=west, fill=white, inner sep=0.8pt] at (46,-4) {オレンジ：まっすぐの道};
\end{tikzpicture}
\bkfigcaption{主役の山を登る2つの道。重力は保存力なので、どちらの道でも、重力に逆らってする仕事は $mg\times(100-83)$ で同じ。遠回りすると、坂がゆるやかになって必要な力は小さくなるが、進む距離が長くなる。}
\label{bk:8-7:fig-yama}
\end{bkfigure}

\kviiihhead{\bkhatenaref の答え}
\bkhatenaref の力 $\vec F_2=(-y,\ x)$ では、仕事は道によって変わります（道A 1\,J、道B $-1$\,J、ななめ 0\,J）。$\vec F_2$ は保存力ではないからです。いっぽう、重力や $\vec F_1$ のような保存力では、遠回りしても仕事は変わりません。

どんな力が保存力かは、平面の力なら $\dfrac{\partial F_y}{\partial x}-\dfrac{\partial F_x}{\partial y}$ が0かどうかで見分けられます（すきまのない広がりの中なら。8-10 の rot）。$\vec F_1$ は $1-1=0$、$\vec F_2$ は $1-(-1)=2$ です。

\kviiihhead{場合 H：閉じた道 ∮}
はじめの点とおわりの点が同じで、1周してもどる道を\textbf{閉じた道}といい、線積分を次のように書きます（\textbf{周回積分}。しゅうかいせきぶん）。\bkidx{しゅうかいせきぶん}{周回積分（∮）}
\kviiihline{$\displaystyle\oint_C\vec F\cdot\dd\vec r$}
∮ の◯は「道が閉じている（1周してもとにもどる）」という印です。円の形の道という意味ではありません。正方形でも三角形でも、閉じていれば ∮ です。どの道かは、右下の $C$ が決めます。
\begin{itemize}[leftmargin=1.2\zw, itemsep=2pt, topsep=0pt, label={\color{\bkTT}●}]
\item 保存力なら、はじめとおわりが同じなので、$U(\text{はじめ})-U(\text{はじめ})$ で、$\oint_C\vec F\cdot\dd\vec r=0$。
\item 保存力でなければ、0とはかぎりません。$\vec F_2$ で、1辺1の正方形を左回りに1周すると2、半径1の円では $2\pi$ です。どちらも「2×囲んだ面積」になっています。このことは、8-11 のストークスの定理（グリーンの定理）で、「rot の値×面積」として学びます。
\end{itemize}

\bksubsection{10}{場合 K：状態の図の上の道}

\label{bk:8-7:pv}線積分の「道」は、地図の上の道だけではありません。シリンダーの中の気体の、体積 $V$ を横、圧力 $p$ をたてにとった図（$p$-$V$ 図）を考えます。気体の状態は、この図の上の点で表され、状態が変わると点が動いて道をえがきます。気体が外にする仕事は、5-12 と同じく $\int p\,\dd V$ で、これも道にそってたす積分です。

\begin{bkfigure}
\begin{tikzpicture}[x=1mm, y=1mm]
  \draw[bkInk, line width=0.7pt, -{Stealth[length=2mm]}] (0,0) -- (78,0) node[right, font=\footnotesize, inner sep=1pt] {$V$（L）};
  \draw[bkInk, line width=0.7pt, -{Stealth[length=2mm]}] (0,0) -- (0,52) node[above, font=\footnotesize, inner sep=1pt] {$p$（kPa）};
  \foreach \kviiihx in {1,2,3} {\draw[bkInk, line width=0.5pt] ({\kviiihx*22},0) -- ({\kviiihx*22},-1.2); \node[font=\footnotesize, anchor=north] at ({\kviiihx*22},-1.5) {\kviiihx};}
  \foreach \kviiihy in {100,200} {\draw[bkInk, line width=0.5pt] (0,{\kviiihy*0.2}) -- (-1.2,{\kviiihy*0.2}); \node[font=\footnotesize, anchor=east] at (-1.5,{\kviiihy*0.2}) {\kviiihy};}
  \fill[bkOrange!15] (22,0) -- (22,40) -- (66,40) -- (66,0) -- cycle;
  \draw[bkOrange!85!black, line width=1.4pt, -{Stealth[length=2.4mm]}] (22,40) -- (66,40) -- (66,20);
  \draw[bkBlue, line width=1.4pt, dash pattern=on 3pt off 1.5pt, -{Stealth[length=2.4mm]}] (22,40) -- (22,20) -- (66,20);
  \draw[bkPurple!80!black, line width=1.2pt, dash pattern=on 1pt off 1.2pt] (22,40) -- (66,20);
  \fill[bkInk] (22,40) circle (0.9) (66,20) circle (0.9);
  \node[font=\footnotesize, anchor=south] at (44,41) {道1：400\,J};
  \node[font=\footnotesize, anchor=north, fill=white, inner sep=0.8pt] at (44,18.5) {道2：200\,J};
  \node[font=\footnotesize, anchor=west, fill=white, inner sep=0.8pt] at (46,33) {まっすぐ：300\,J};
\end{tikzpicture}
\bkfigcaption{$p$-$V$ 図の上の3つの道。どれも 1\,L・200\,kPa の状態から 3\,L・100\,kPa の状態へ。気体がする仕事 $\int p\,\dd V$ は、道の下の面積で、道によって変わる（例題11）。}
\label{bk:8-7:fig-pv}
\end{bkfigure}

道1は、圧力 200\,kPa のまま 1\,L から 3\,L にふくらみ、そのあと体積を変えずに圧力を下げます。体積が変わらない区間は $\dd V=0$ なので、仕事は0。だから道1の仕事は $200\times10^3\times2\times10^{-3}=400$\,J です（場合 A と同じく、区間ごとに動く文字は1つ）。道2は先に圧力を下げるので 200\,J、まっすぐの道は 300\,J です。気体がする仕事は、同じ2つの状態を結んでも道によって変わります。物理の熱力学（熱-2）では、このことから「仕事や熱は道によって変わるが、内部エネルギーは状態だけで決まる」ことを学びます。

\begin{seikaku}
\kviiihbox{道の表し方を変えても、答えは同じ}
同じ道を、別の数で表してもかまいません。たとえば、ななめの道（$(0,\ 0)$ から $(1,\ 1)$）を、$\vec r(u)=(u^2,\ u^2)$（$0\leqq u\leqq1$）と表すと、$\vec r\,'(u)=(2u,\ 2u)$ です。$\vec F_1=(y,\ x)$ では
\kviiihline{$\displaystyle\int_0^1(u^2,\ u^2)\cdot(2u,\ 2u)\,\dd u=\int_0^1 4u^3\,\dd u=1$}
で、$\vec r(t)=(t,\ t)$ で計算した1と同じです。歩く速さを変えても、1歩の矢印 $\vec r\,'\,\dd t$ の中で、速さと時間の変化が打ち消し合うからです（5-16 の置換積分）。ただし、はじめとおわりを入れかえると、符号が変わります。
\end{seikaku}

\begin{motto}[5]
\kviiihbox{道を選んで、積分をいちばん小さくする（変分法の入口）}
この章では、道を決めてから線積分を計算しました。反対に、「線積分がいちばん小さくなる道はどれか」を求める問題もあります。たとえば、2点を結ぶ道の長さ $\int_C\dd s$ がいちばん小さいのは、まっすぐな線分です。両はしを持ってつり下げた鎖（長さは決まっている）は、位置エネルギー（$\int_C\lambda gy\,\dd s$）がいちばん小さくなる形になり、その形は 6-11 で出てきた懸垂線です。光は、かかる時間がいちばん短い道を進みます（フェルマーの原理）。このような問題をあつかう数学を\textbf{変分法}（へんぶんほう）といい、そのための式を\textbf{オイラー・ラグランジュ方程式}といいます。物理の本（波-3）で使います。
\end{motto}

\begin{senmonsho}[線積分の専門書ふうの書き方]
\begin{bkhonbun}
曲線 $C:\ \symbf{r}=\symbf{r}(t)\ (a\le t\le b)$ に沿うベクトル場 $\symbf{F}$ の線積分を $\displaystyle\int_C\symbf{F}\cdot d\symbf{r}=\int_a^b\symbf{F}(\symbf{r}(t))\cdot\frac{d\symbf{r}}{dt}\,dt$ で定める．平面では $\displaystyle\int_C P\,dx+Q\,dy$ とも書く．$\symbf{F}=\nabla\phi$ ならば $\displaystyle\int_C\symbf{F}\cdot d\symbf{r}=\phi(\symbf{r}(b))-\phi(\symbf{r}(a))$．
\end{bkhonbun}
\bkyaku{太字の $\symbf{F}$・$\symbf{r}$ は、この本の矢印の $\vec F$・$\vec r$。「沿う」は「そって」。$P\,dx+Q\,dy$ は $F_x\,\dd x+F_y\,\dd y$ のこと（$P=F_x$、$Q=F_y$）。数学の本では、$\symbf{F}=\nabla\phi$（$+$ の符号）と書くことが多く、そのとき答えは「おわり－はじめ」になる。物理の本は $\vec F=-\mathrm{grad}\,U$（$-$ の符号）なので「はじめ－おわり」。符号のつけ方がちがうだけで、同じこと。物理の本では $\displaystyle\oint\vec E\cdot\dd\vec l$ のように $\dd\vec l$ や $\dd\vec s$ と書くこともある。}
\end{senmonsho}

\begin{taisetsu}
\kviiihbox{線積分の3つの段}
\kviiihdan{1}\enspace 記号を読む：何を（$\vec F\cdot\dd\vec r$ か $f\,\dd s$）・どの道で（$C$ と向き）\par
\kviiihdan{3}\enspace 直す前に：データ・数値？　$p$-$V$ 図？　$U$ のわかる保存力？　閉じた道？　力が式でない？\par
\kviiihdan{2}\enspace 道を1つの数で表し、$\dd\vec r=\vec r\,'(t)\,\dd t$ で、ふつうの定積分1回\par
\[ \kviiihN{\int_C\vec F\cdot\dd\vec r}=\kviiihS{\int_a^b\vec F(\vec r(t))\cdot\vec r\,'(t)\,\dd t} \]
\kviiihbox{場合分けの表8-7（まとめ）}
\begin{bktable}{colspec={Q[c,m,9mm] Q[l,m,44mm] X[l,m]}, row{2-Z}={font=\small}, rowsep=2pt, colsep=3pt}
記号 & こんなとき & 直し方 \\
\kviiihcase{A} & 軸に平行 & $F_x\,\dd x$ か $F_y\,\dd y$（もう1つは止める） \\
\kviiihcase{B} & まっすぐなななめ & $\vec A+t(\vec B-\vec A)$、$0\leqq t\leqq1$ \\
\kviiihcase{C} & $y=g(x)$ & $(1,\ g'(x))\,\dd x$、$y$ に $g(x)$ を入れる \\
\kviiihcase{D} & 円・円弧 & $(-R\sin\theta,\ R\cos\theta)\,\dd\theta$ \\
\kviiihcase{E} & 折れ線 & 区間ごとにたす。逆向きは符号が変わる \\
\kviiihcase{F} & 力が式で書いてない & 0 か $\pm F\times$長さ \\
\kviiihcase{G} & 保存力 & $U(\text{はじめ})-U(\text{おわり})$ \\
\kviiihcase{H} & 閉じた道 ∮ & 保存力なら0 \\
\kviiihcase{I} & 長さでたす & $\dd s=|\vec r\,'|\,\dd t$。向きによらない \\
\kviiihcase{J} & データだけ・数値で & 表計算で $\sum\vec F_i\cdot\Delta\vec r_i$ \\
\kviiihcase{K} & $p$-$V$ 図 & 区間ごとに $\int p\,\dd V$ \\
\end{bktable}
\end{taisetsu}

\begin{kotoba}
\small
\bkword{線積分}{道を小さな1歩に分け、「力と1歩の矢印の内積」（または「量×1歩の長さ」）をたし、細かくした先。操作の名前。}
\bkword{閉じた道}{はじめとおわりが同じ道。線積分は ∮ で書く。◯は「閉じている」という印。}
\bkword{保存力}{$\vec F=-\mathrm{grad}\,U$ と書ける力。仕事は道によらず、$U(\text{はじめ})-U(\text{おわり})$。閉じた道では0。}
\bkword{$\dd s$}{道にそった1歩の長さ。$\dd s=|\vec r\,'(t)|\,\dd t$。}
\bkword{$p$-$V$ 図}{気体の体積を横、圧力をたてにとった図。仕事 $\int p\,\dd V$ は道で変わる。}
\end{kotoba}

\bkkigocards{\bkkigocard{\scalebox{0.62}{$\displaystyle\int_C\!\vec F\cdot\dd\vec r$}}{インテグラル・シー、エフ・ドット・ディーアール}{道 $C$ にそって、力と1歩の矢印の内積をたした先（線積分）。道には向きがある}{line integral}{8-7}}
{\bkkigocard{$\displaystyle\oint_C$}{周回積分（しゅうかいせきぶん）}{閉じた道 $C$ を1周する線積分。◯は「閉じている」印で、円の形の意味ではない}{closed line integral (circulation)}{8-7}}




%% ---------------- 例題 ----------------
\bksection{reidai}{例題でたしかめよう}

\noindent 例題の「考え方」の1行めには、場合分けの表8-7のどの場合かを書きました。力の単位は N、長さの単位は m、仕事の単位は J です。

\kviiihhead{例題を、流れ図でたどってみよう}
例題1〜11を、図\ref{bk:8-7:fig-flow}・図\ref{bk:8-7:fig-flow2}の流れ図の上から順にたどると、次のようになります（書いていない問いは「いいえ」。力が2つある例題は $\vec F_2$ で）。どれも、考え方の1行めの場合に着きます。じぶんでとこうも同じ道すじです。
\begin{bkcenter}
\begin{bktable}{colspec={Q[c,m,10mm] Q[l,m,28mm] X[l,m] Q[c,m,16mm]}, row{2-Z}={font=\small}, rowsep=0.8pt, colsep=3pt}
例題 & 題 & 「はい」と答えた問い & 着いた所 \\
1 & 軸に平行な道 & 軸に平行な線分だけ？ & \kviiihcase{A} \\
2 & ななめの道 & まっすぐなななめ？ & \kviiihcase{B} \\
3 & 放物線の道 & $y=g(x)$ のグラフ？ & \kviiihcase{C} \\
4 & 円を1周 & 閉じた道？ → 円・円弧？ & \kviiihcase{H}\kviiihcase{D} \\
5 & 三角形の道 & 閉じた道？ → ななめ・曲線のある折れ線？（区間ごとに A・A・B） & \kviiihcase{H}\kviiihcase{E} \\
6 & 摩擦力 & 力が式で書いてない？ & \kviiihcase{F} \\
7 & 保存力 & $U$ のわかる保存力？ & \kviiihcase{G} \\
8 & 正方形の1周 & 閉じた道？ → 軸に平行な線分だけ？ & \kviiihcase{H}\kviiihcase{A} \\
9 & 半円の針金 & 名前の段で $f\,\dd s$（I）→ 円・円弧？ & \kviiihcase{I}\kviiihcase{D} \\
10 & 表計算で円 & データだけ・数値で？ & \kviiihcase{J} \\
11 & $p$-$V$ 図 & $p$-$V$ 図の上の道？ & \kviiihcase{K} \\
\end{bktable}
\bktabcaption{例題1〜11の流れ図の道すじ。}
\end{bkcenter}
{\small 例題1〜5・8 の $\vec F_1$ は保存力ですが、問題に $U$ が書いていないので、道ごとに計算して「道によらない」ことを確かめます（あとから G で確かめられる）。}

例題を読む前に、表の「はい」と答えた問いを、図\ref{bk:8-7:fig-flow}・図\ref{bk:8-7:fig-flow2}の上で指でたどってみましょう。考え方の1行めの「場合」と同じ札に着けば、その例題で使う直し方が先にわかります。

\begin{reidai}[軸に平行な道]
力 $\vec F_2=(-y,\ x)$ と $\vec F_1=(y,\ x)$ について、道A（$(0,\ 0)\to(1,\ 0)\to(1,\ 1)$）と道B（$(0,\ 0)\to(0,\ 1)\to(1,\ 1)$）の線積分を、それぞれ求めましょう。\par
\kangaekata \kviiihbaai{道が横とたての線分}{8-7}{A}\par
\kviiihdan{1}\enspace $\int_C\vec F\cdot\dd\vec r$：道にそって、力と1歩の内積をたす。\par
\kviiihdan{2}\enspace 横の区間は $F_x\,\dd x$（$y$ を止める）、たての区間は $F_y\,\dd y$（$x$ を止める）。\par
\shiki\kviiihs\hangindent=2.6\zw\hangafter=1 $\vec F_2$、道A：$\int_0^1(-0)\,\dd x+\int_0^1 1\,\dd y=0+1=1$\par
\kviiihb （右へは $y=0$ で $F_x=0$。上へは $x=1$ で $F_y=1$）\par
\kviiiha $\vec F_2$、道B：$\int_0^1 0\,\dd y+\int_0^1(-1)\,\dd x=0-1=-1$\par
\kviiihb （上へは $x=0$ で $F_y=0$。右へは $y=1$ で $F_x=-1$）\par
\kviiiha $\vec F_1$、道A：$\int_0^1 0\,\dd x+\int_0^1 1\,\dd y=1$\quad （$F_x=y=0$、$F_y=x=1$）\par
\kviiiha $\vec F_1$、道B：$\int_0^1 0\,\dd y+\int_0^1 1\,\dd x=1$\quad （$F_y=x=0$、$F_x=y=1$）\par
\kotae $\vec F_2$：道A 1\,J、道B $-1$\,J。$\vec F_1$：どちらも 1\,J
\end{reidai}

\begin{jibun}
力 $\vec F_3=(2x,\ 2y)$ と $\vec F_4=(y,\ -x)$ について、道Aと道Bの線積分を、それぞれ求めましょう。
\bkjans{$\vec F_3$：道A $\int_0^1 2x\,\dd x+\int_0^1 2y\,\dd y=1+1=2$、道B も $1+1=2$。$\vec F_4$：道A $\int_0^1 0\,\dd x+\int_0^1(-1)\,\dd y=-1$、道B $\int_0^1 0\,\dd y+\int_0^1 1\,\dd x=1$。}
\end{jibun}

\begin{reidai}[まっすぐなななめの道]
原点 $(0,\ 0)$ から点 $(1,\ 1)$ へまっすぐ進む道で、次の2つの力の線積分を求めましょう。\par
\kviiihline{$\vec F_1=(y,\ x)$、\quad $\vec F_2=(-y,\ x)$}
\kangaekata \kviiihbaai{まっすぐなななめの道}{8-7}{B}\par
\kviiihdan{1}\enspace ななめの道にそって、力と1歩の内積をたす。\par
\kviiihdan{2}\enspace $\vec r(t)=(0,\ 0)+t(1,\ 1)=(t,\ t)$、$0\leqq t\leqq1$、$\vec r\,'(t)=(1,\ 1)$。\par
\begin{center}\begin{tikzpicture}[x=0.75mm, y=0.75mm]
  \draw[bkLine, line width=0.4pt, step=20] (0,0) grid (40,40);
  \draw[bkPurple!80!black, line width=1.4pt, -{Stealth[length=2.2mm]}] (0,0) -- (40,40);
  \foreach \kviiiht in {0.25,0.5,0.75,1} {\draw[bkGray!80!black, line width=0.8pt, -{Stealth[length=1.6mm]}] ({40*\kviiiht},{40*\kviiiht}) -- ++({-9*\kviiiht},{9*\kviiiht});
    \draw[bkOrange!85!black, line width=0.8pt, -{Stealth[length=1.6mm]}] ({40*\kviiiht},{40*\kviiiht}) -- ++({9*\kviiiht},{9*\kviiiht});}
  \node[font=\footnotesize, anchor=west, align=left] at (56,30) {灰色：$\vec F_2=(-y,\ x)$（道と直角）};
  \node[font=\footnotesize, anchor=west, align=left] at (56,20) {オレンジ：$\vec F_1=(y,\ x)$（道と同じ向き）};
  \node[font=\footnotesize, anchor=north] at (0,-1.5) {$(0,\ 0)$};
  \node[font=\footnotesize, anchor=north] at (40,-1.5) {$(1,\ 0)$};
\end{tikzpicture}\end{center}
\shiki\kviiihs\hangindent=2.6\zw\hangafter=1 $\vec F_1(\vec r(t))=(t,\ t)$、内積 $(t,\ t)\cdot(1,\ 1)=2t$\par
\kviiiha $\int_0^1 2t\,\dd t=\Bigl[\,t^2\,\Bigr]_0^1=1$\par
\kviiiha $\vec F_2(\vec r(t))=(-t,\ t)$、内積 $(-t,\ t)\cdot(1,\ 1)=0$\par
\kviiiha $\int_0^1 0\,\dd t=0$\par
\kotae $\vec F_1$：1\,J、$\vec F_2$：0\,J（$\vec F_2$ は、この道の上ではいつも道と直角）
\end{reidai}

\begin{jibun}
点 $(1,\ 0)$ から点 $(0,\ 1)$ へまっすぐ進む道で、$\vec F_1=(y,\ x)$ と $\vec F_2=(-y,\ x)$ の線積分を求めましょう。
\bkjans{$\vec r(t)=(1-t,\ t)$、$\vec r\,'=(-1,\ 1)$。$\vec F_1=(t,\ 1-t)$ で内積 $-t+1-t=1-2t$、$\int_0^1(1-2t)\,\dd t=0$。$\vec F_2=(-t,\ 1-t)$ で内積 $t+1-t=1$、$\int_0^1 1\,\dd t=1$。}
\end{jibun}

\begin{reidai}[放物線の道]
原点から点 $(1,\ 1)$ へ、放物線 $y=x^2$ にそって進む道で、$\vec F_1=(y,\ x)$ と $\vec F_2=(-y,\ x)$ の線積分を求めましょう。\par
\kangaekata \kviiihbaai{道が $y=g(x)$ のグラフ}{8-7}{C}\par
\kviiihdan{1}\enspace 放物線にそって、力と1歩の内積をたす。\par
\kviiihdan{2}\enspace $x$ を道の数にする：$\vec r=(x,\ x^2)$、$\dd\vec r=(1,\ 2x)\,\dd x$、$0\leqq x\leqq1$。$y$ には $x^2$ を入れる。\par
\shiki\kviiihs\hangindent=2.6\zw\hangafter=1 $\vec F_1=(x^2,\ x)$、内積 $x^2\times1+x\times2x=3x^2$\par
\kviiiha $\int_0^1 3x^2\,\dd x=\Bigl[\,x^3\,\Bigr]_0^1=1$\par
\kviiiha $\vec F_2=(-x^2,\ x)$、内積 $-x^2+x\times2x=x^2$\par
\kviiiha $\int_0^1 x^2\,\dd x=\Bigl[\,\frac{1}{3}x^3\,\Bigr]_0^1=\frac{1}{3}$\par
\kotae $\vec F_1$：1\,J（ほかの道と同じ）、$\vec F_2$：$\dfrac{1}{3}$\,J
\end{reidai}

\begin{jibun}
原点から点 $(1,\ 1)$ へ、曲線 $y=x^3$ にそって進む道で、$\vec F_1=(y,\ x)$ と $\vec F_2=(-y,\ x)$ の線積分を求めましょう。
\bkjans{$\dd\vec r=(1,\ 3x^2)\,\dd x$。$\vec F_1=(x^3,\ x)$ で内積 $x^3+3x^3=4x^3$、$\int_0^1 4x^3\,\dd x=1$。$\vec F_2=(-x^3,\ x)$ で内積 $-x^3+3x^3=2x^3$、$\int_0^1 2x^3\,\dd x=\frac{1}{2}$。}
\end{jibun}

\begin{reidai}[円を1周する道]
中心が原点、半径3の円を、点 $(3,\ 0)$ から左回りに1周する道で、$\vec F_2=(-y,\ x)$ と $\vec F_1=(y,\ x)$ の線積分を求めましょう。\par
\kangaekata \kviiihbaai{閉じた道で、道が円}{8-7}{H,D}\par
\kviiihdan{1}\enspace 円にそって1周、力と1歩の内積をたす（閉じた道なので ∮ と書く）。\par
\kviiihdan{2}\enspace $\vec r=(3\cos\theta,\ 3\sin\theta)$、$\dd\vec r=(-3\sin\theta,\ 3\cos\theta)\,\dd\theta$、$0\leqq\theta\leqq2\pi$。\par
\shiki\kviiihs\hangindent=2.6\zw\hangafter=1 $\vec F_2=(-3\sin\theta,\ 3\cos\theta)$、内積 $9\sin^2\theta+9\cos^2\theta=9$\par
\kviiiha $\int_0^{2\pi}9\,\dd\theta=18\pi\approx56.5$\par
\kviiiha $\vec F_1=(3\sin\theta,\ 3\cos\theta)$、内積 $-9\sin^2\theta+9\cos^2\theta=9\cos2\theta$\par
\kviiiha $\int_0^{2\pi}9\cos2\theta\,\dd\theta=\Bigl[\,\frac{9}{2}\sin2\theta\,\Bigr]_0^{2\pi}=0$\par
\kotae $\vec F_2$：$18\pi$（$=2\times$円の面積 $9\pi$）、$\vec F_1$：0（保存力なので1周で0）
\end{reidai}

\begin{jibun}
中心が原点、半径2の円を左回りに1周する道で、$\vec F_2=(-y,\ x)$ と $\vec F_1=(y,\ x)$ の線積分を求めましょう。
\bkjans{$\dd\vec r=(-2\sin\theta,\ 2\cos\theta)\,\dd\theta$。$\vec F_2$ の内積は4で、$\int_0^{2\pi}4\,\dd\theta=8\pi$。$\vec F_1$ の内積は $4\cos2\theta$ で、$\int_0^{2\pi}4\cos2\theta\,\dd\theta=0$。}
\end{jibun}

\begin{reidai}[折れ線の道（三角形）]
道Aで $(0,\ 0)\to(1,\ 0)\to(1,\ 1)$ と進み、ななめの線分を通って $(0,\ 0)$ へもどる三角形の道で、$\vec F_2=(-y,\ x)$ と $\vec F_1=(y,\ x)$ の線積分を求めましょう。\par
\kangaekata \kviiihbaai{閉じた道で、ななめの区間もある折れ線}{8-7}{H,E}\par
\kviiihdan{1}\enspace 3つの区間をつないだ閉じた道。区間ごとに直してたす。\par
\kviiihdan{2}\enspace 横・たての区間は A、ななめの区間は B。もどる区間は例題2の道の逆向きなので、符号を変える。\par
\shiki\kviiihs\hangindent=2.6\zw\hangafter=1 $\vec F_2$：道Aの部分は例題1より1。もどる区間は例題2の0の逆で0\par
\kviiihb 合計 $1+0=1$\par
\kviiiha $\vec F_1$：道Aの部分は1。もどる区間は例題2の1の逆で $-1$\par
\kviiihb 合計 $1+(-1)=0$\par
\kotae $\vec F_2$：1、$\vec F_1$：0
\end{reidai}

\begin{jibun}
道Bで $(0,\ 0)\to(0,\ 1)\to(1,\ 1)$ と進み、ななめの線分を通って $(0,\ 0)$ へもどる三角形の道で、$\vec F_2=(-y,\ x)$ と $\vec F_1=(y,\ x)$ の線積分を求めましょう。
\bkjans{$\vec F_2$：道Bの部分 $-1$、もどる区間 0（ななめの0の逆）、合計 $-1$。$\vec F_1$：道Bの部分 1、もどる区間 $-1$、合計 0。}
\end{jibun}

\begin{reidai}[摩擦力の仕事]
質量 \qty{10}{kg} の箱を、床の上で点Pから点Qまで、すべらせて動かします。動摩擦係数は $\mu'=0.3$ です。（1）PからQまでまっすぐ \qty{5}{m} 動かすとき、（2）横に \qty{3}{m}、たてに \qty{4}{m} の L字の道で動かすとき、動摩擦力がする仕事をそれぞれ求めましょう（重力加速度 \qty{9.8}{m/s^2}）。\par
\kangaekata \kviiihbaai{力が式で書いてなく、いつも道と反対向きで大きさ一定}{8-7}{F}\par
\kviiihdan{1}\enspace 道にそって、摩擦力と1歩の内積をたす。\par
\kviiihdan{3}\enspace 摩擦力はいつも動く向きと反対で大きさ一定。内積は $-$（力の大きさ）×（1歩の長さ）。\par
\shiki\kviiihs\hangindent=2.6\zw\hangafter=1 床が箱をおし返す力（垂直抗力）$N$〔N〕は $10\times9.8=98$\par
\kviiiha 動摩擦力の大きさ $\mu'N=0.3\times98=\qty{29.4}{N}$\par
\kviiiha （1）$-29.4\times5=-147$\par
\kviiiha （2）道の長さ $3+4=\qty{7}{m}$、$-29.4\times7=-205.8$\par
\kotae （1）$-147$\,J　（2）約 $-206$\,J（遠回りすると、摩擦の仕事は大きく負になる）
\end{reidai}

\begin{jibun}
質量 \qty{5}{kg} の箱を、床の上で、（1）直径 \qty{4}{m} の半円の道、（2）その直径にそったまっすぐな道で動かします。動摩擦係数 $\mu'=0.4$ のとき、動摩擦力がする仕事をそれぞれ求めましょう（重力加速度 \qty{9.8}{m/s^2}）。
\bkjans{動摩擦力 $0.4\times5\times9.8=\qty{19.6}{N}$。（1）半円の長さ $2\pi\approx\qty{6.28}{m}$、仕事 $-19.6\times2\pi\approx-123$\,J。（2）$-19.6\times4=-78.4$\,J。}
\end{jibun}

\begin{reidai}[保存力の仕事]
（1）質量 \qty{50}{kg} の人が、主役の山 $h=100-0.01x^2-0.02y^2$〔m〕の点 $(30,\ 20)$ から頂上までゆっくり登ります。重力がする仕事と、人がする仕事を求めましょう（重力加速度 \qty{9.8}{m/s^2}）。\par
（2）$\vec F_1=(y,\ x)$（$U=-xy$）が、点 $(1,\ 2)$ から点 $(3,\ 1)$ まで、どんな道で動かしたときもする仕事を求めましょう。\par
\kangaekata \kviiihbaai{$U$ のわかる保存力}{8-7}{G}\par
\kviiihdan{1}\enspace 道にそって、力と1歩の内積をたす（道はわからなくてよい）。\par
\kviiihdan{3}\enspace 保存力なので $U(\text{はじめ})-U(\text{おわり})$。重力は $U=mgz$、$\vec F_1$ は $U=-xy$。\par
\shiki\kviiihs\hangindent=2.6\zw\hangafter=1 （1）はじめの高さ $h(30,\ 20)=100-9-8=\qty{83}{m}$、おわり \qty{100}{m}\par
\kviiihb 重力の仕事 $=mg\times83-mg\times100=50\times9.8\times(83-100)=-8330$\par
\kviiihb 人の仕事は、重力に逆らうので $+8330$\par
\kviiiha （2）$U(1,\ 2)-U(3,\ 1)=(-1\times2)-(-3\times1)=-2+3=1$\par
\kotae （1）重力 $-8330$\,J、人 8330\,J（どの道でも同じ）　（2）1\,J
\end{reidai}

\begin{jibun}
（1）質量 \qty{2}{kg} の玉が、高さ \qty{5}{m} の所から、曲がった坂をころがって地面まで下ります。重力がする仕事を求めましょう（重力加速度 \qty{9.8}{m/s^2}）。\par
（2）$\vec F_1=(y,\ x)$（$U=-xy$）が、点 $(0,\ 2)$ から点 $(2,\ 3)$ まで動かしたときにする仕事を求めましょう。
\bkjans{（1）$U=mgz$ で、$2\times9.8\times5-0=98$\,J（坂の形によらない）。（2）$U=-xy$ で、$U(0,\ 2)-U(2,\ 3)=0-(-6)=6$\,J。}
\end{jibun}

\begin{reidai}[閉じた道の線積分]
1辺1の正方形の道 $(0,\ 0)\to(1,\ 0)\to(1,\ 1)\to(0,\ 1)\to(0,\ 0)$（左回り）で、$\oint\vec F_2\cdot\dd\vec r$ と $\oint\vec F_1\cdot\dd\vec r$ を求めましょう（$\vec F_2=(-y,\ x)$、$\vec F_1=(y,\ x)$）。\par
\kangaekata \kviiihbaai{閉じた道で、軸に平行な線分だけ}{8-7}{H,A}\par
\kviiihdan{1}\enspace 正方形を1周して、力と1歩の内積をたす。\par
\kviiihdan{3}\enspace $\vec F_1$ は保存力なので、1周で0。$\vec F_2$ は保存力でないので、区間に分ける（A）。\par
\kviiihdan{2}\enspace 下の辺（右へ、$y=0$）、右の辺（上へ、$x=1$）、上の辺（左へ、$y=1$）、左の辺（下へ、$x=0$）。\par
\shiki\kviiihs\hangindent=2.6\zw\hangafter=1 下の辺：$\int_0^1 0\,\dd x=0$\quad （$F_x=-y=0$）\par
\kviiiha 右の辺：$\int_0^1 1\,\dd y=1$\quad （$F_y=x=1$）\par
\kviiiha 上の辺：$\int_1^0(-1)\,\dd x=1$\quad （$F_x=-1$、左へ進むので1から0まで）\par
\kviiiha 左の辺：$\int_1^0 0\,\dd y=0$\quad （$F_y=x=0$）\par
\kviiiha 合計 $0+1+1+0=2$\par
\kotae $\oint\vec F_2\cdot\dd\vec r=2$（$=2\times$面積1）、$\oint\vec F_1\cdot\dd\vec r=0$
\end{reidai}

\begin{jibun}
1辺2の正方形の道 $(0,\ 0)\to(2,\ 0)\to(2,\ 2)\to(0,\ 2)\to(0,\ 0)$ で、$\oint\vec F_2\cdot\dd\vec r$ と $\oint\vec F_1\cdot\dd\vec r$ を求めましょう。
\bkjans{$\vec F_2$：下の辺0、右の辺 $\int_0^2 2\,\dd y=4$、上の辺 $\int_2^0(-2)\,\dd x=4$、左の辺0、合計8（$=2\times$面積4）。$\vec F_1$：保存力なので0。}
\end{jibun}

\begin{reidai}[半円の針金の質量]
半径 \qty{0.5}{m} の半円の形の針金（$\theta$ が0から $\pi$）があります。角 $\theta$ の所の1\,m あたりの質量は $\lambda=2+\sin\theta$〔kg/m〕です。針金の質量を求めましょう。\par
\kangaekata \kviiihbaai{長さでたす（向きのない量）で、道が円}{8-7}{I,D}\par
\kviiihdan{1}\enspace $\int_C\lambda\,\dd s$：針金にそって「$\lambda$×1歩の長さ」をたす。\par
\kviiihdan{2}\enspace 円なので $\dd s=R\,\dd\theta=0.5\,\dd\theta$、$0\leqq\theta\leqq\pi$。\par
\begin{center}\begin{tikzpicture}[x=1mm, y=1mm]
  \foreach \kviiiha in {0,6,...,174} {\pgfmathsetmacro\kviiihw{0.6+1.2*sin(\kviiiha+3)}\draw[bkOrange!85!black, line width=\kviiihw mm] (\kviiiha:20) arc (\kviiiha:{\kviiiha+6.3}:20);}
  \draw[bkInk, line width=0.6pt, -{Stealth[length=2mm]}] (-25,0) -- (27,0) node[right, font=\footnotesize, inner sep=1pt] {$x$};
  \fill[bkInk] (0,0) circle (0.6);
  \draw[bkInk, line width=0.5pt] (0,0) -- (60:20);
  \draw[bkInk, line width=0.5pt] (6,0) arc (0:60:6);
  \node[font=\footnotesize] at (30:8.5) {$\theta$};
  \node[font=\footnotesize, anchor=west, align=left] at (30,12) {太いほど重い：$\lambda=2+\sin\theta$};
  \node[font=\footnotesize, anchor=west, align=left] at (30,5) {まん中（$\theta=90°$）が3\,kg/m、両はしが2\,kg/m};
\end{tikzpicture}\end{center}
\shiki\kviiihs\hangindent=2.6\zw\hangafter=1 $\int_0^{\pi}(2+\sin\theta)\times0.5\,\dd\theta=0.5\Bigl[\,2\theta-\cos\theta\,\Bigr]_0^{\pi}$\par
\kviiiha $=0.5\times\bigl((2\pi+1)-(0-1)\bigr)=0.5\times(2\pi+2)=\pi+1$\par
\kotae $\pi+1\approx\qty{4.14}{kg}$
\end{reidai}

\begin{jibun}
半径 \qty{2}{m} の4分の1円の形の針金（$\theta$ が0から $\frac{\pi}{2}$）で、1\,m あたりの質量が $\lambda=1+\cos\theta$〔kg/m〕のとき、針金の質量を求めましょう。
\bkjans{$\dd s=2\,\dd\theta$。$\int_0^{\pi/2}(1+\cos\theta)\times2\,\dd\theta=2\Bigl[\theta+\sin\theta\Bigr]_0^{\pi/2}=2\Bigl(\frac{\pi}{2}+1\Bigr)=\pi+2\approx\qty{5.14}{kg}$。}
\end{jibun}

\begin{reidai}[表計算で円の道をたす]
半径1の円を、点 $(1,\ 0)$、$(0,\ 1)$、$(-1,\ 0)$、$(0,\ -1)$ を順に結ぶ4本の線分（正方形）で近づけ、$\vec F_2=(-y,\ x)$ の線積分の近い値を、線分のまん中の点の力を使って求めましょう。\par
\kangaekata \kviiihbaai{数値で近い値を求める}{8-7}{J}\par
\kviiihdan{1}\enspace 円の道にそって、力と1歩の内積をたす。\par
\kviiihdan{3}\enspace 4本の線分に分け、「まん中の点の力 $\vec F_i$」と「線分の矢印 $\Delta\vec r_i$」の内積をたす。\par
\shiki\kviiihs\hangindent=2.6\zw\hangafter=1 1本め：$\Delta\vec r=(0,\ 1)-(1,\ 0)=(-1,\ 1)$、まん中 $(0.5,\ 0.5)$\par
\kviiihb $\vec F_2=(-0.5,\ 0.5)$、内積 $0.5+0.5=1$\par
\kviiiha 2本め：$\Delta\vec r=(-1,\ -1)$、まん中 $(-0.5,\ 0.5)$、$\vec F_2=(-0.5,\ -0.5)$\par
\kviiihb 内積 $0.5+0.5=1$\par
\kviiiha 3本め・4本めも、同じように1ずつ（円の対称性）\par
\kviiiha 合計 $1\times4=4$\par
\kotae 4（本当の値 $2\pi\approx6.28$ より小さい。表\ref{bk:8-7:tab-n}のように、線分をふやすと近づく）
\end{reidai}

\begin{jibun}
半径1の円を、角 $0°$、$60°$、$120°$、…、$300°$ の6つの点を順に結ぶ6本の線分で近づけ、$\vec F_2=(-y,\ x)$ の線積分の近い値を、線分のまん中の点の力を使って求めましょう。
\bkjans{1本め：$(1,\ 0)$ から $(0.5,\ 0.866)$、$\Delta\vec r=(-0.5,\ 0.866)$、まん中 $(0.75,\ 0.433)$、$\vec F_2=(-0.433,\ 0.75)$、内積 $0.2165+0.6495=0.866$。6本とも同じなので、$0.866\times6\approx5.20$。}
\end{jibun}

\begin{reidai}[$p$-$V$ 図の上の仕事]
シリンダーの中の気体が、体積 \qty{1}{L}・圧力 \qty{200}{kPa} の状態から、体積 \qty{3}{L}・圧力 \qty{100}{kPa} の状態へ変わります。気体がする仕事 $\int p\,\dd V$ を、次の3つの道で求めましょう。（1）道1：圧力 \qty{200}{kPa} のままふくらみ、そのあと体積を変えずに圧力を下げる。（2）道2：先に体積を変えずに \qty{100}{kPa} まで下げ、そのあとふくらむ。（3）$p$-$V$ 図の上をまっすぐ進む。\par
\kangaekata \kviiihbaai{状態の図の上の道}{8-7}{K}\par
\kviiihdan{1}\enspace $p$-$V$ 図の道にそって「$p\times\dd V$」をたす。\par
\kviiihdan{2}\enspace 区間ごとに：体積が変わらない区間は $\dd V=0$ で0。圧力が一定の区間は $p\times$（体積の変化）。\par
\shiki\kviiihs\hangindent=2.6\zw\hangafter=1 体積の変化 \qty{2}{L} $=2\times10^{-3}$\,m$^3$（\qty{1000}{L/m^3}）\par
\kviiiha （1）$200\times10^3\times2\times10^{-3}+0=400$\par
\kviiiha （2）$0+100\times10^3\times2\times10^{-3}=200$\par
\kviiiha （3）道の上で $p=200-50(V-1)$〔kPa、$V$ は L〕\par
\kviiihb $\displaystyle\int_1^3\bigl(200-50(V-1)\bigr)\dd V=\Bigl[\,200V-25(V-1)^2\,\Bigr]_1^3$\par
\kviiihb $=(600-100)-200=300$（kPa・L）\par
\kviiihb \qty{1}{kPa}×\qty{1}{L} $=10^3\times10^{-3}=\qty{1}{J}$ なので 300\,J（下の台形の面積 $\dfrac{200+100}{2}\times2$ と同じ）\par
\kotae （1）400\,J　（2）200\,J　（3）300\,J（道によって変わる）
\end{reidai}

\begin{jibun}
気体が、体積 \qty{2}{L}・圧力 \qty{300}{kPa} の状態から、体積 \qty{5}{L}・圧力 \qty{100}{kPa} の状態へ変わります。例題11と同じ3つの道で、気体がする仕事を求めましょう。
\bkjans{体積の変化 $3\times10^{-3}$\,m$^3$。道1：$300\times10^3\times3\times10^{-3}=900$\,J。道2：$100\times10^3\times3\times10^{-3}=300$\,J。まっすぐ：平均の圧力 \qty{200}{kPa} で、$200\times10^3\times3\times10^{-3}=600$\,J。}
\end{jibun}

%% ---------------- まちがえやすいところ ----------------
\bksection{tsumazuki}{まちがえやすいところをチェック}

\begin{tsumazuki}{① 向きを忘れて、力の大きさ×長さにする}
\bkng{道B での $\vec F_2$ の仕事＝$|\vec F_2|\times$（道の長さ）で正の値}
\bkok{$\vec F\cdot\dd\vec r$（内積）をたす。道Bでは $-1$}
\[ \kviiihN{\int_C\vec F_2\cdot\dd\vec r}=\kviiihS{\int_0^1 0\,\dd y+\int_0^1(-1)\,\dd x}=-1 \]
仕事は内積なので、力と進む向きが反対なら負、直角なら0です。$|\vec F|\,\dd s$ をたしてよいのは、力がいつも道と同じ向きのとき（場合 F）だけです。
\end{tsumazuki}

\begin{tsumazuki}{② $t$ の範囲を、はじめとおわりで逆にする}
\bkng{$(1,\ 1)$ から $(0,\ 0)$ へ進むのに、$\vec r=(t,\ t)$ で $\int_0^1$ とする}
\bkok{$\vec r=(1-t,\ 1-t)$ で $\int_0^1$、または $\vec r=(t,\ t)$ で $\int_1^0$}
$t$ の下のはしは「はじめの点」、上のはしは「おわりの点」です。逆にすると符号が逆になります。
\end{tsumazuki}



\begingroup\tcbset{equal height group=kviiihEHT}%
\noindent\begin{minipage}[t]{0.485\linewidth}\vspace{0pt}
\begin{tsumazuki}{③ ∮ の◯を「円の道」と思う}
\raggedright\bkng{$\oint$ だから、道は円}
\bkok{◯は「閉じた道」という印。正方形でも三角形でもよい}
どの道かは、右下の $C$ と、問題の文で決まります。8-8 の $\oiint$（閉じた面）の◯も同じ意味です。
\end{tsumazuki}
\end{minipage}\hfill
\begin{minipage}[t]{0.485\linewidth}\vspace{0pt}
\begin{tsumazuki}{④ 保存力でない力に、$U$ の差を使う}
\raggedright\bkng{$\vec F_2$ の仕事は、はじめとおわりだけで決まる}
\bkok{$\vec F_2$ は保存力ではない。道ごとに計算する}
$U$ の差が使えるのは、$\vec F=-\mathrm{grad}\,U$ と書ける力だけです。摩擦力も保存力ではありません。
\end{tsumazuki}
\end{minipage}\par
\endgroup

\begingroup\tcbset{equal height group=kviiihEHU}%
\noindent\begin{minipage}[t]{0.485\linewidth}\vspace{0pt}
\begin{tsumazuki}{⑤ $y=g(x)$ の道で、$y$ を残したまま積分する}
\raggedright\bkng{$y=x^2$ の道で、\\$\int_0^1(-y+2x^2)\,\dd x=-y+\frac{2}{3}$}
\bkok{$y$ に $x^2$ を入れて、\\$\int_0^1(-x^2+2x^2)\,\dd x=\frac{1}{3}$}
道の上では $y$ は $x$ で決まるので、積分する前に $y$ を $g(x)$ に置きかえます。
\end{tsumazuki}
\end{minipage}\hfill
\begin{minipage}[t]{0.485\linewidth}\vspace{0pt}
\begin{tsumazuki}{⑥ 左や下へ進む区間を $\int_0^1$ にする}
\raggedright\bkng{正方形の上の辺（$(1,\ 1)$ から $(0,\ 1)$）を $\int_0^1(-1)\,\dd x=-1$}
\bkok{$\int_1^0(-1)\,\dd x=1$}
左へ進むとき、$x$ は1から0へ。下のはしがはじめの $x$ です。
\end{tsumazuki}
\end{minipage}\par
\endgroup





%% ---------------- 物理・工学ではこう使う ----------------
\begin{butsuri}
\kviiihbox{エネルギーの保存：ジェットコースター}
\noindent\begin{minipage}[c]{0.66\linewidth}
\raggedright
重力は保存力なので、ジェットコースターがどんな曲がったレールを通っても、重力の仕事は高さの差だけで決まります。だから、まさつが小さければ、レールの形を細かく計算しなくても、高さから速さがわかります（物理 力-4）。実際には、まさつや空気の力（保存力でない）の分、少しずつエネルギーが失われます。
\end{minipage}\hfill
\begin{minipage}[c]{0.3\linewidth}
\centering
\begin{tikzpicture}[x=0.85mm, y=1mm]
  \draw[bkInk, line width=1.1pt] plot[smooth] coordinates {(0,22) (8,6) (16,14) (24,3) (34,10)};
  \filldraw[fill=bkOrange!60, draw=bkInk, line width=0.6pt] (1.5,18) rectangle ++(4,3);
  \draw[bkInk, line width=0.5pt, dash pattern=on 1pt off 1pt] (0,22) -- (36,22) (24,3) -- (36,3);
  \draw[bkInk, line width=0.6pt, {Stealth[length=1.6mm]}-{Stealth[length=1.6mm]}] (35,3) -- (35,22);
  \node[font=\footnotesize, anchor=west] at (35.5,12.5) {$h$};
\end{tikzpicture}
\end{minipage}
\par\vspace{4pt}
\kviiihbox{電圧と起電力}
電場 $\vec E$ の中で、電気をもった粒が動くときも、仕事は線積分です。電圧（電位差）は、電場を道にそって線積分したものです。止まった電荷が作る電場は保存力の場なので、電圧は道によらず決まります。いっぽう、発電機のように磁場が時間とともに変わる所では、閉じた道の線積分 $\oint\vec E\cdot\dd\vec r$ が0でなくなり、これが\textbf{起電力}です（物理 電-8、8-14）。
\par\vspace{4pt}
\kviiihbox{電流をはさんではかる：クランプメーター}
\noindent\begin{minipage}[c]{0.66\linewidth}
\raggedright
電線を流れる電流のまわりには、輪の形の磁場ができます。磁場を電線を囲む閉じた道で線積分すると、囲んだ電流に比例します（アンペールの法則、8-11）。電線のまわりの円の道では、磁場は道と同じ向きで大きさが一定なので、場合 F と同じく「磁場の大きさ×円周の長さ」です。クランプメーターは、電線をはさむ輪で、この磁場をはかり、電線を切らずに電流を知る計測器です。
\end{minipage}\hfill
\begin{minipage}[c]{0.3\linewidth}
\centering
\begin{tikzpicture}[x=1mm, y=1mm]
  \filldraw[fill=bkGray!50, draw=bkInk, line width=0.6pt] (0,0) circle (2);
  \fill[bkInk] (0,0) circle (0.6);
  \foreach \kviiihr in {6,10} {\draw[bkBlue, line width=0.9pt, -{Stealth[length=1.8mm]}] (\kviiihr,0) arc (0:340:\kviiihr);}
  \node[font=\footnotesize, anchor=north, align=center] at (0,-12) {電流と、\\まわりの磁場};
\end{tikzpicture}
\end{minipage}
\end{butsuri}

\begin{column}{エンジンの1回の仕事は、閉じた道の線積分}
自動車のエンジンのシリンダーの中では、気体が「吸う・おしちぢめる・燃えてふくらむ・出す」をくり返します。$p$-$V$ 図の上では、気体の状態は1回ごとに同じ所へもどる閉じた道をえがきます。気体が1回にする仕事は $\oint p\,\dd V$ で、これは道が囲む面積です。この面積が大きいほど、1回にたくさん仕事をするエンジンです（物理 熱-6）。
\end{column}

\kviiihhead{流れ図でたどってみよう（まとめの練習）}
章末問題に進む前に、図\ref{bk:8-7:fig-flow}・図\ref{bk:8-7:fig-flow2}の流れ図を、もう一度上からたどってみましょう。次の線積分は、表8-7のどの場合でしょう。
\begin{enumerate}[label={\bkcircled[\bkTT]{\arabic*}}, leftmargin=2.4\zw, itemsep=3pt, topsep=2pt]
\item $(0,\ 0)$ から $(2,\ 4)$ まで、放物線 $y=x^2$ にそって、$\vec F=(y,\ 1)$\par\nobreak\noindent 答え：\kviiihcase{C}。保存力かどうかは書いていない（$U$ がない）。閉じていない道で、$y=g(x)$ のグラフ。
\item 気体を、圧力 \qty{100}{kPa} のまま、\qty{2}{L} から \qty{5}{L} までふくらませる\par\nobreak\noindent 答え：\kviiihcase{K}。$p$-$V$ 図の上の、横にまっすぐな道。
\item 重さ \qty{20}{N} の荷物を、床から高さ \qty{1.5}{m} のたなまで、いろいろな道で持ち上げるときの重力の仕事（$U=20z$）\par\nobreak\noindent 答え：\kviiihcase{G}。$U$ のわかる保存力なので、道によらず $0-30=\qty{-30}{J}$。
\item $(0,\ 0)\to(3,\ 0)\to(3,\ 2)$ の折れ線にそって、$\vec F=(x+y,\ xy)$\par\nobreak\noindent 答え：\kviiihcase{A}。ななめ・曲線の区間がなく、軸に平行な線分だけ。
\item 長方形 $(0,\ 0)\to(2,\ 0)\to(2,\ 1)\to(0,\ 1)\to(0,\ 0)$ を1周、$\vec F=(-y,\ x)$\par\nobreak\noindent 答え：\kviiihcase{H}→\kviiihcase{A}。閉じた道なので H、1周を下の問いでたどると、軸に平行な線分だけ。
\item 1秒ごとに記録した位置と、そのときの風の力のデータから、風がした仕事を求める\par\nobreak\noindent 答え：\kviiihcase{J}。式がなく、データだけ。
\end{enumerate}

%% ---------------- 章末問題 ----------------
\bkendproblems
\bkproblemgroup{基本}{5}

\bkproblem 力 $\vec F=(2y,\ x)$ について、道A（$(0,\ 0)\to(1,\ 0)\to(1,\ 1)$）と道B（$(0,\ 0)\to(0,\ 1)\to(1,\ 1)$）の線積分を求めましょう\mbox{（\bkreidairef{1}）}。
\bkpans{道A：右へは $F_x=2y=0$ で0、上へは $F_y=x=1$ で1、合計1。道B：上へは $F_y=0$ で0、右へは $F_x=2$ で2、合計2。}

\bkproblem 原点から点 $(1,\ 1)$ へまっすぐ進む道で、$\vec F=(2y,\ x)$ の線積分を求めましょう\mbox{（\bkreidairef{2}）}。
\bkpans{$\vec r=(t,\ t)$、$\vec r\,'=(1,\ 1)$。$\vec F=(2t,\ t)$ で内積 $3t$。$\int_0^1 3t\,\dd t=1.5$。}

\bkproblem 原点から点 $(1,\ 1)$ へ放物線 $y=x^2$ にそって進む道で、$\vec F=(2y,\ x)$ の線積分を求めましょう\mbox{（\bkreidairef{3}）}。
\bkpans{$\dd\vec r=(1,\ 2x)\,\dd x$。$\vec F=(2x^2,\ x)$ で内積 $2x^2+2x^2=4x^2$。$\int_0^1 4x^2\,\dd x=\frac{4}{3}$。}

\bkproblem 中心が原点、半径2の円の上半分を、点 $(2,\ 0)$ から点 $(-2,\ 0)$ まで左回りに進む道で、$\vec F_2=(-y,\ x)$ の線積分を求めましょう\mbox{（\bkreidairef{4}）}。
\bkpans{$\vec r=(2\cos\theta,\ 2\sin\theta)$、$0\leqq\theta\leqq\pi$。内積は4。$\int_0^{\pi}4\,\dd\theta=4\pi$。}

\bkproblem 半径 \qty{1}{m} の4分の1円の形の針金（$\theta$ が0から $\frac{\pi}{2}$）で、1\,m あたりの質量が $\lambda=3+\sin\theta$〔kg/m〕のとき、針金の質量を求めましょう\mbox{（\bkreidairef{9}）}。
\bkpans{$\dd s=\dd\theta$。$\int_0^{\pi/2}(3+\sin\theta)\,\dd\theta=\Bigl[3\theta-\cos\theta\Bigr]_0^{\pi/2}=\frac{3\pi}{2}+1\approx\qty{5.71}{kg}$。}

\bkproblem 質量 \qty{60}{kg} の人が、主役の山の点 $(40,\ 10)$ から頂上までゆっくり登ります。人がする仕事を求めましょう（重力加速度 \qty{9.8}{m/s^2}）\mbox{（\bkreidairef{7}）}。
\bkpans{$h(40,\ 10)=100-16-2=\qty{82}{m}$。$60\times9.8\times(100-82)=\num{10584}$、約 \qty{1.06e4}{J}（どの道でも同じ）。}

\bkproblemgroup{チャレンジ}{5}

\bkproblem 道Aで $(0,\ 0)\to(1,\ 0)\to(1,\ 1)$ と進み、ななめの線分を通って $(0,\ 0)$ へもどる三角形の道で、$\vec F=(2y,\ x)$ の線積分を求めましょう\mbox{（\bkreidairef{5}）}。
\bkpans{道Aの部分は1（問題1）。もどる区間は問題2の1.5の逆で $-1.5$。合計 $-0.5$。}

\bkproblem 1辺1の正方形の道 $(0,\ 0)\to(1,\ 0)\to(1,\ 1)\to(0,\ 1)\to(0,\ 0)$ で、$\oint\vec F\cdot\dd\vec r$ を求めましょう（$\vec F=(2y,\ x)$）\mbox{（\bkreidairef{8}）}。
\bkpans{下の辺0、右の辺 $\int_0^1 1\,\dd y=1$、上の辺 $\int_1^0 2\,\dd x=-2$、左の辺0。合計 $-1$。}

\bkproblem 質量 \qty{20}{kg} の箱を、動摩擦係数 $\mu'=0.25$ の床の上で、（1）まっすぐ \qty{10}{m}、（2）横に \qty{6}{m}、たてに \qty{8}{m} の L字の道で動かすとき、動摩擦力がする仕事を求めましょう（重力加速度 \qty{9.8}{m/s^2}）\mbox{（\bkreidairef{6}）}。
\bkpans{動摩擦力 $0.25\times20\times9.8=\qty{49}{N}$。（1）$-49\times10=-490$\,J。（2）$-49\times14=-686$\,J。}

\bkproblem 気体が、体積 \qty{1}{L}・圧力 \qty{400}{kPa} の状態から、体積 \qty{4}{L}・圧力 \qty{100}{kPa} の状態へ変わります。例題11と同じ3つの道で、気体がする仕事を求めましょう\mbox{（\bkreidairef{11}）}。
\bkpans{体積の変化 $3\times10^{-3}$\,m$^3$。道1：$400\times10^3\times3\times10^{-3}=1200$\,J。道2：$100\times10^3\times3\times10^{-3}=300$\,J。まっすぐ：$250\times10^3\times3\times10^{-3}=750$\,J。}

\begin{furikaeri}
\begin{checklist}
\item 道を小さな1歩に分け、力と1歩の内積をたして仕事を表計算で求め、線積分の記号を読める。\bkfurref{bk:8-7:ippo}
\item 軸に平行な道（折れ線）の線積分を、ふつうの定積分に直して計算できる。\bkfurref{bk:8-7:jiku}
\item 場合分けの表を使って、ななめ・曲線・円の道、保存力、閉じた道の線積分を計算できる。\bkfurref{bk:8-7:basho}
\end{checklist}
\end{furikaeri}

\bkfukidashi{mina}{道にそってたすときは、道を1つの数で表せばふつうの定積分になるんだね。次は、面を通りぬける量をたす面積分。}
